% Section 5 of the paper = v3 §4 (G3, the queue).
\section{The queue}\label{sec:queue}
% v3: §4 (opening paragraph)

This section proves that, on an event of probability greater than $\tfrac12$ (Lemma~\ref{lem:budget}), the Lindley queue of
Section~\ref{sec:construction} satisfies $Q_j\le B(s_j)$ at every slot (Theorem~\ref{thm:queue}). Here $B(x)$ is quasi-polylogarithmic in
$x$, and is bounded by $H(x)^{O(1)}$ on the raised range of Section~\ref{sec:queue-capl}. Consequently (L-a) holds:
$|D(x)|\le B(x)+\tfrac{H(x)}2(1+\kappa_+/\log x)+3\le x^{\eta}$. The ingredients are:
% Approved deviation from v3 (CORRECTIONS-vs-v3.md, #15): v3 has H(x)/2+2.
\begin{itemize}
\item the landing count and the local count (Lemma~\ref{lem:landing} and Corollary~\ref{cor:localcount});
\item the good event and the first-failure budget (Definition~\ref{def:good} and Lemma~\ref{lem:budget});
\item the cluster event $\qEcl$ and its per-prime increment bound (Proposition~\ref{prop:ecl} and Corollary~\ref{cor:increment});
\item the chain of lemmas proving $\qEcl$ at every scale (Sections~\ref{sec:queue-ecl}, \ref{sec:queue-cover} and~\ref{sec:queue-capl}, together
  with Appendix~\ref{app:rankfull});
\item the analysis of the absorbed set $\Abs$: its pointwise sparsity (Lemma~\ref{lem:lemmaP}) and the consequences
  (Section~\ref{sec:queue-abs}).
\end{itemize}

\subsection{Conventions, parameters and standing facts}\label{sec:queue-conv}\label{sec:queue-standing}
% v3: §4.1

\paragraph{Scales.}
A \emph{dyadic scale} is $x=2^mX_0$ with $m\ge0$. We write $L:=\log x$ and $\ell:=\log L$. All parameters are \emph{frozen per
scale}: they are evaluated at $x$ and used for heights in $[x,2x)$. We write $\Lambda:=\log X_0\ge10^8$
(Section~\ref{sec:params}).

\paragraph{Parameters at scale $x$.}
These are the choices of Section~\ref{sec:params} and Section~\ref{sec:construction}.
\begin{itemize}
\item $\delta_1:=2L^{-6}$. The \emph{grid windows} at height $z\in[x,2x]$ are the intervals $[k\delta_1,(k+2)\delta_1)$,
  $k\in\Z$, of length $2\delta_1=4L^{-6}$.
\item $r_0:=\lceil L/(3\ell)\rceil+3$ is the rank threshold (Theorem~\ref{thm:R}).
\item $c:=c(x)$ is the capacity threshold, fixed in Theorem~\ref{thm:caplw}. We put $Q:=c^{1/2}$ and
  $\varphi(u):=c+uQ+u(u+1)/2$.
\item $M_1:=r_0+\varphi(r_0)+1$ is the cluster bound. Then $M_1\le 2c$, and $M_1\ge L$ for $L\ge L_*$.
\item $H:=H(x):=\max(\qHstd(x),M_1(x)^2)$, where $\qHstd(x):=\exp(8\ell^2)$. This is the width of the box jitter of
  Section~\ref{sec:construction}.
  \begin{itemize}
  \item Where $c(x)=\tfrac14\qHstd(x)^{1/2}$ we have $M_1\le\tfrac12\qHstd^{1/2}$, so $H=\qHstd$.
  \item The \emph{raised range} is the set of scales where $c(x)>\tfrac14\qHstd(x)^{1/2}$, that is, where
    $\gamma(s)\Lambda>4\ell^2-\log4$ (Theorem~\ref{thm:caplw}). There $H=\max(\qHstd,M_1^2)\ge\qHstd$.
  \item In all cases the tolerance $M_1\le H/L^{14}$ of Section~\ref{sec:params} holds.
  \end{itemize}
\item $L_0:=\max(L^5,48M_1L^3/c'^2)$ is the block length. We put $q_{\min}:=L_0L^6$ and $y:=H\cdot L_0\cdot L^6$. The latter is
  the absorption threshold of Section~\ref{sec:construction}, including its ramp.
\item $B:=B(x)$ is the queue bound of Theorem~\ref{thm:queue}. The constants $c'=0.15$ and $K_d=3$ are the drift and edge constants of
  the drift bound R1$'$ (Corollary~\ref{cor:r1}).
\item $\lambda:=4L$ in Theorem~\ref{thm:queue}. For the budget (Lemma~\ref{lem:budget}) we use, per scale $u$,
  $\lambda_u:=3\log u+2\log\log u$ and $\lambda'_u:=\lambda_u+2+\log H(u)$.
\end{itemize}
In summary,
\begin{equation}\label{eq:params}
\begin{gathered}
r_0=\Big\lceil\frac{L}{3\ell}\Big\rceil+3,\qquad M_1=r_0+\varphi(r_0)+1,\qquad H=\max(\qHstd,M_1^2),\\
L_0=\max\Big(L^5,\frac{48M_1L^3}{c'^2}\Big),\qquad y=HL_0L^6.
\end{gathered}
\end{equation}

\paragraph{Filtration.}
Slots $s_1<s_2<\dots$ (Section~\ref{sec:construction}) are processed in increasing order. $\mathcal F_t$ is the $\sigma$-field of
all decisions and jitter offsets at slots $\le t$, and $\mathcal F_{s-}$ is the same for slots $<s$. For a decision slot $s$,
$\NN_{\prec s}$ denotes the set of elements of $\NN$ all of whose primes are decided at slots $<s$. It is
$\mathcal F_{s-}$-measurable.

\paragraph{Convention on $X_0$.}
With centred placement (Section~\ref{sec:construction}), jittered primes have values at least
$\Xm:=X_0-H(X_0)/2=X_0(1-O(e^{2715-\Lambda}))$. Wherever this section or Appendix~\ref{app:rankfull} says that ``a jittered prime is
at least $X_0$'' (in (F3$'$), (V1), Lemma~\ref{lem:rowweights}, the $\tau_j$ of Appendix~\ref{app:rankfull}, and $\mathrm{Proj}_t$), read
$\Xm$. Every such use has a margin far larger than $e^{-10^7}$, and Lemma~\ref{lem:rowweights}(e) already uses the weaker bound
$\tfrac23(X_0+0.99(i-1))$.

\paragraph{The Lindley recursion.}
By Lemma~\ref{lem:lindley}, the queue is $Q_0=0$ and
\begin{equation}\label{eq:lindley}
Q_j=\max(0,\,Q_{j-1}+A_j-1),
\end{equation}
where $A_j$ is the number of arrivals assigned to slot $s_j$. The \emph{excess} of a slot window $W$ is
$\qexc(W):=C(W)-\#\{\text{slots of }W\}$. By the excess identity (Lemma~\ref{lem:excess}),
$\qexc(W)=\qCjit(W)-\qljit(W)\pm2$. Here $\qCjit$ counts the non-absorbed composites whose latest-decided prime is
jittered, and $\qljit(W)$ counts the jittered (continuous) slots of $W$. Lindley's formula gives
\begin{equation}\label{eq:lindleymax}
Q_j=\max\Big(0,\ \max_{k<j}\qexc((k,j])\Big).
\end{equation}

\paragraph{Standing facts at height $z$.}
All of the following are sure (pathwise) statements, except (F2), which is distributional.
\begin{enumerate}[label=\textup{(F\arabic*)},leftmargin=*]
\item\label{eq:sF1} At most $H(t)(1+\kappa_+/\log t)+3$ decided primes have jitter windows containing a given point $t$, since the slot
  density satisfies $T'\le1+\kappa_+/\log t$ (Section~\ref{sec:construction}). Consequently, the label sums over slots in
  Lemma~\ref{lem:pil} carry the factor $\theta:=1+1/\Lambda$, which bounds $1+\kappa_+/\log q+3/H$ for jittered primes
  $q\ge\Xm$ (with $\kappa_+=0.8$).
  % Approved deviation from v3 (CORRECTIONS-vs-v3.md, #12): v3 has theta := 1+1/L.
\item\label{eq:sF2} Given $\mathcal F_{s-}$, the position of the prime decided at slot $s$ has density at most $1/H(s)$ (centred box
  jitter). The slots of distinct primes are independent.
\end{enumerate}
\begin{enumerate}[label=\textup{(F3$'$)},leftmargin=*]
\item\label{eq:sF3p} \emph{(Harmonic sums, $\mathcal F_{s-}$-form.)} On $\Good(<s)$ (Definition~\ref{def:good}), for every $Y\ge1$,
\[
\sum_{c\in\NN_{\prec s},\ c\le Y}\frac1c\ \le\ A_s(Y):=
\begin{cases}
\log Y+2, & \text{if } Y+H(Y)\le s,\\[2pt]
e^{\gamma}\log s\cdot\exp\big(2.21/\Lambda+5H(s)/s\big), & \text{otherwise.}
\end{cases}
\]
Moreover $A_s(Y)\le A':=1.8(\log z+1)$ for $s\le 2z$ and $Y\le 2z$.
% Approved deviation from v3 (CORRECTIONS-vs-v3.md, #2): v3 has 5H(s)/X_0 in place of 5H(s)/s.
\end{enumerate}
\begin{proof}
If $Y+H(Y)\le s$, then all of $\NN\cap[1,Y]$ is decided before $s$, and the bound is the sure harmonic bound of
Lemma~\ref{lem:cm}(iv), which holds on $\Good(<s)$. Otherwise we bound the sum by
$\prod_{p\text{ decided before }s}(1-1/p)^{-1}$.
\begin{itemize}
\item The rational primes $p<X_0$ give at most $e^{\gamma}\Lambda(1+1/(2\Lambda^2))$ \cite[Thm.~8, (3.29)]{RosserSchoenfeld1962}, valid for
  $X_0\ge286$.
\item The realised primes in $[X_0,s-H(s)]$ give at most $(\log s/\Lambda)e^{2.20/\Lambda+1/X_0}$, by the sure local Mertens
  bound of Lemma~\ref{lem:cm}(iv). That bound reads $\NN$ only on $[1,s-H]$ and holds on $\Good(<s)$.
\item The remaining primes decided before $s$ have values in $(s-H,s+H/2]$, hence values greater than $s/2$. There are at most
  $2H+1$ of them, so they contribute a factor at most $(1-2/s)^{-(2H+1)}\le e^{5H(s)/s}$. Since
  $\log H\le8\ell^2+10^{-6}\log s$ (Theorem~\ref{thm:caplw}(c)), this factor is at most $e^{10^{-9}}$. \qedhere
\end{itemize}
\end{proof}
\begin{enumerate}[label=\textup{(F4)},leftmargin=*]
\item\label{eq:sF4} \emph{(Sure Mertens and Chebyshev bounds on $\Good(<s)$.)} These are the realised-prime bounds of
  Lemma~\ref{lem:cm}(iv), used inside (F3$'$) and in Appendix~\ref{app:rankfull}.
\end{enumerate}

\paragraph{Absorbed and non-absorbed elements.}
$\Abs$ (see \eqref{eq:absorbed}) is the extended absorption set: the elements $n$ with $\log n\ge L_*=10^7$ whose latest jittered
prime $r$ satisfies $r<y(n)$. In Lemma~\ref{lem:lemmaP}, $y(n)$ is replaced by $\sup_{[u,2u]}y$ on the dyadic block of
$n$. Absorbed elements are counted in the target $T$ and are not arrivals. All statements about arrivals concern
$\NN\setminus\Abs$.
\begin{itemize}
\item \emph{Stopped process.} Every statement of this section is proved on $\Good(<t)$ for the relevant $t$, that is, for the
  process stopped at the first failure time $\tau$. There Lemma~\ref{lem:lemmaP} gives $\rho_{\Abs}\le L^{-4}<0.1/\log u$, so the
  capped density $\bar\rho_{\Abs}=\min(\rho_{\Abs},0.1/\log u)$ of \eqref{eq:target} equals $\rho_{\Abs}$. No argument of this
  section uses $\rho_{\Abs}$ off the good event: Lemma~\ref{lem:lemmaP}, Proposition~\ref{prop:S1S2}, the compensator of
  $\sigma_{\Abs}$ (Lemma~\ref{lem:rhoA}(b)) and \eqref{eq:T1} all run on the good event.
\item \emph{Slot lengths.} This section uses $\rhodet$ only through $T'\le1+1/(\kappa u)\le1.01$, in the linear
  normalisation of Section~\ref{sec:construction} (see \eqref{eq:rhodet} and Lemma~\ref{lem:blin}(1)): $\rhodet\le$ (rigid fraction)
  $+1/(\kappa u)$, and the positivity margin is \eqref{eq:G1}, with no $\kappa/2$ loss. This enters as slot spacing at least
  $0.99$ in (F1) and Appendix~\ref{app:rankfull}, and in the candidate counts of Lemma~\ref{lem:landing}. No $\kappa/2$ form is used.
\end{itemize}

\subsection{Landing count and local count}\label{sec:queue-landing}
% v3: §4.2

Fix a dyadic $P$ and put $\lambda_P:=4\log P$, $H_-:=H(P/2)$, $H_+:=H(8P)$ and $\rho_H:=H_+/H_-$. Then $\rho_H\le1.35$ for
$\log P\ge1100$, and $\rho_H\le1+1.2\cdot10^{-5}$ for $\log P\ge10^8$, the raise of Theorem~\ref{thm:caplw} included, because
$d\log H/d\log P\le16\ell/L+10^{-7}$ (Theorem~\ref{thm:caplw}(c)). For a window $J\subset[P,4P+1]$ with $|J|\le6H(P)$ we
define:
\begin{itemize}
\item the \emph{candidates} $\mathcal C(J)$: the decision slots $s$ with $|s-J|\le H(s)/2$. (Selection is by decision slot,
  which is predictable; selection by the value $q$ would not be.)
\item $\Lambda(J)$: the number of primes decided at slots in $\mathcal C(J)$ that land in $J$;
\item the compensator $\mu(J):=\sum_{s\in\mathcal C(J)}\Prob(\text{the prime at $s$ lands in $J$}\mid\mathcal F_{s-})$;
\item $\bar\mu(J):=(|J|+H_++2)\min(1,|J|/H_-)$. By (F1)--(F2) this is a \emph{sure} bound:
  $\mu(J)\le\bar\mu(J)\le(1+\rho_H)|J|+2$.
\end{itemize}

\begin{lemma}[Landing count]\label{lem:landing}
% v3: Lemma 4.1
Let $\qLand(J)$ be the event
\begin{equation}\label{eq:landing}
|\Lambda(J)-\mu(J)|\le\sqrt{2\bar\mu(J)\lambda_P}+\tfrac23\lambda_P,\qquad\text{hence}\qquad
\Lambda(J)\le\min\big(\#\mathcal C(J),\,2\bar\mu(J)+\tfrac76\lambda_P\big).
\end{equation}
Then $\Prob(\qLand(J)^c)\le2e^{-\lambda_P}=2P^{-4}$, with no conditioning. Summed over the landing windows $J\in\mathcal J_P$
defined below, the total is at most $146H(P)P^{-3}$. The event $\qLand(J)$ is $\mathcal F_t$-measurable with
$t(\qLand(J)):=\max J+H(\max J)$.
\end{lemma}

\begin{proof}
Take the indicators $\mathbf 1(\text{the prime at $s$ lands in $J$})$, $s\in\mathcal C(J)$, in slot order. Their conditional
means sum to $\mu(J)\le\bar\mu(J)$ surely, and the increments are at most $1$. Apply Freedman's inequality
\cite[(4.1)]{Freedman1975} to both tails. We use it in the form
\begin{equation}\label{eq:freedman}
\Prob(\exists n:\ S_n\ge a,\ V_n\le v)\le\exp\Big(-\frac{a^2}{2v+2ab/3}\Big),
\end{equation}
for a martingale $S_n$ with increments at most $b$ and predictable quadratic variation $V_n$, which follows from \cite[(4.1)]{Freedman1975} by scaling and Bennett's bound $(1+u)\log(1+u)-u\ge u^2/(2(1+u/3))$. Take
$v=\bar\mu$ and $b=1$. The second form in \eqref{eq:landing} follows from $\sqrt{2\bar\mu\lambda}\le\bar\mu+\lambda/2$.
\end{proof}

\begin{remark}\label{rem:oldlanding}
% v3: §4.2, Remark after Lemma 4.1
A bound of the form $\Lambda(J)\le|J|(1+|J|/H)+\lambda$ with failure probability $e^{-\lambda}$ would need $\lambda\ge3.92\mu$,
which is false for $\lambda\lesssim|J|$; it would also condition on later decisions. It is not used.
\end{remark}

\paragraph{Landing windows.} Let $\mathcal J_P$ be the set of windows $[P+i/2,\,P+j/2)$ with integers $0\le i<j$,
$j-i\le12H(P)$ and $P+j/2\le4P+1$. There are at most $(6P+3)\cdot12H(P)\le73P\,H(P)$ of them. Every $J\subset[P,4P]$ with
$|J|\le6H(P)-1$ is contained in a smallest $J^+\in\mathcal J_P$, and $|J^+|\le|J|+1$.
% Approved deviation from v3 (CORRECTIONS-vs-v3.md, #20): v3 leaves the window family of the union bound unspecified ("72P H(P) grid windows", total 144H(P)P^{-3}).

\begin{corollary}[Local count]\label{cor:localcount}
% v3: Corollary 4.2
Let $J\subset[P,4P]$ with $|J|\le6H(P)$ (including $|J|<1$). On the queue bound at the slots $s_j\le\max J+H(\max J)$, and,
if $|J|\le6H(P)-1$, on $\qLand(J^+)$,
\begin{equation}\label{eq:localcount}
|N(J)-T(J)|\le2B(4P)+2\Lambda(J)+2|J|+5\le(4\rho_H+6)|J|+\qKloc(P)\le11.4|J|+\qKloc(P),
\end{equation}
where $\qKloc(P):=2B(4P)+\tfrac73\lambda_P+23$.
% Approved deviation from v3 (CORRECTIONS-vs-v3.md, #20): v3 has +13 in K_loc; the extra 10 pays for the passage from J to the landing window J^+.
Under the target \eqref{eq:target} of Section~\ref{sec:construction},
$T(J)\le2|J|+2+\#\Abs(J)$, so the non-absorbed count satisfies $N(J)-\#\Abs(J)\le13.4|J|+\qKloc(P)+2$.
\end{corollary}

\begin{proof}
The count identity (Lemma~\ref{lem:count}), applied to the slots meeting $J$, gives $|\Delta Q|\le B$, $|\Delta C|\le B+1$,
$|\Delta(j-T)|\le2$ and $|\Delta L_\pm(J)|\le\Lambda(J)+\#\{\text{slots in }J\}\le\Lambda(J)+|J|+1$. This gives the first
inequality. Now let $|J|\le6H(P)-1$. Since $J\subset J^+$, we have $\mathcal C(J)\subset\mathcal C(J^+)$ and
$\Lambda(J)\le\Lambda(J^+)$. On $\qLand(J^+)$, by \eqref{eq:landing},
\[
2\Lambda(J)\le4\bar\mu(J^+)+\tfrac73\lambda_P\le4(1+\rho_H)(|J|+1)+8+\tfrac73\lambda_P\le4(1+\rho_H)|J|+\tfrac73\lambda_P+17.4,
\]
using $\rho_H\le1.35$; this gives the second inequality. If $6H(P)-1<|J|\le6H(P)$, then $|J|\ge H$ and Remark~\ref{rem:longwindows}
gives the bound directly.
\end{proof}

\begin{remark}\label{rem:longwindows}
% v3: §4.2, after Corollary 4.2
Windows longer than $6H(P)$ need no landing event. The sure lag bound \eqref{eq:lag} gives
$|N(J)-T(J)|\le2\sup|D|\le2B+H(1+\kappa_+/L)+6\le1.01|J|+2B+6$ for $|J|\ge H$.
% Approved deviation from v3 (CORRECTIONS-vs-v3.md, #15): v3 has 2B+H+4 <= |J|+2B+4.
Alternatively, splitting $J$ into $\lceil|J|/H\rceil$ pieces gives $N(J)-\#\Abs(J)\le13.5|J|+\qKloc(P)+2$.
\end{remark}

\begin{remark}[Uses of Corollary~\ref{cor:localcount}, and non-circularity]\label{rem:localcount-uses}
% v3: §4.2, "Uses of Corollary 4.2"
Corollary~\ref{cor:localcount} at height $z$ reads the queue bound at slots up to $z+H(z)$ and the event $\qLand$ at height $z$. It is
therefore used only by steps that run after $z+7H(z)$, and always at a strictly smaller height. Its uses are:
\begin{itemize}
\item Lemma~\ref{lem:threshold-mass} (steps (i)--(ii)) and Proposition~\ref{prop:br} (step 3), at heights at most $x/4$ or $\sqrt x$;
\item Lemma~\ref{lem:tau-near}, at heights $p_0\le\sqrt x$;
\item optionally, step (iv) of Lemma~\ref{lem:lemmaP}, at heights $u/d\le u^{1-\theta/2}$.
\end{itemize}
It is \emph{not} used for the per-prime increment bound in Theorem~\ref{thm:queue}. A local count there would force
$B(x)\ge(3/c')L^2B(x/q)$, which is incompatible with \eqref{eq:queueB}. The increment bound comes instead from $\qEcl$
(Corollary~\ref{cor:increment}).
\end{remark}

\subsection{The queue bound}\label{sec:queue-bound}
% v3: §4.3, Theorem 4.3

Fix a dyadic scale $x$ and put $\lambda:=4L$. Recall that $\qljit(\cdot)$ counts continuous (jittered-prime) slots.

\paragraph{Windows.}
\begin{itemize}
\item \emph{Blocks are measured in continuous slots.} $\mathcal W_x$ is the set of slot windows $W=(k,j]$ with right end
  $s_j\in[x,2x)$ and $\qljit(W)\le L_0$, of any total slot count. By \eqref{eq:G1}, continuous slots have density at least
  $0.09/L$, so such a $W$ has physical length at most $(L_0+1)L/0.09\le12L\,L_0$.
\item \emph{Freedman's inequality is applied to deterministic windows.} $\mathcal I_x$ is the set of physical windows
  $I=[a,b)$ with $a,b\in\Z$, $b\in[x,\,2x+12L+2]$ and $b-a\le12L(L_0+1)+2$.
\end{itemize}
For $I\in\mathcal I_x$ and a latest prime $q$, let $X^{(1)}_q(I)$ count the non-absorbed composites $n\in I$ with $q\,\|\,n$
and $q$ latest; let $X^{(2)}_q(I)$ count those with $q^a\,\|\,n$, $a\ge2$; and put $\tilde X_q:=\min(X^{(1)}_q(I),2M_1)$.
These are $\mathcal F_q$-measurable: $n=q^a m$ with $m$ decided before $q$, and absorption is decided by $q<y(n)$. They
depend neither on $T$ nor on the slot positions. Composites have $q\le b/2$, so every $X_q(I)$ is decided by $b/2+H(b/2)$.
The event $\qFr(I)$ below also reads $\qljit(I)$ and uses R1$'$ on $I$, which holds on $\Good(<\max I)$
(Corollary~\ref{cor:r1}); so $\qFr(I)$ carries the time $t(I):=b$.

For $I\in\mathcal I_x$, let $\qFr(I)$ be the event that, with $\ell:=\qljit(I)$ (an integer, possibly random),
\begin{enumerate}[label=\textup{(\alph*)}]
\item $\sum_q\big(\tilde X_q-\E[\tilde X_q\mid\mathcal F_{q-}]\big)\le\sqrt{2v_\ell\lambda}+\tfrac43M_1\lambda$, where
  $v_\ell:=2M_1(\ell+K_d)$;
\item $\#\{q:\ X^{(2)}_q(I)\ge1\}\le\lambda$.
\end{enumerate}
The sums run over latest primes, selected by decision slot.

\begin{theorem}[Queue bound]\label{thm:queue}
% v3: Theorem 4.3
Let $x$ be a dyadic scale. The drift bound R1$'$ is used in the form of Corollary~\ref{cor:r1}, on every window (slot or physical)
in $[1,\infty)$, with $L:=\log(2x+12\log x+2)$; this exceeds $\log x$ by less than $0.7$, which R1$'$'s slack absorbs.
\begin{enumerate}[label=\textup{(\roman*)}]
\item For each $I\in\mathcal I_x$, $\Prob(\qFr(I)^c\cap\Good(<t(I)))\le2(L_0+6)e^{-\lambda}$, with $t(I)=b$. Hence
  $\sum_{I\in\mathcal I_x}\Prob(\qFr(I)^c\cap\Good(<t(I)))\le30L(L_0+1)(L_0+6)x^{-3}\le x^{-2}$.
\item (Deterministic part.) Put $t_j:=\lceil s_{j+1}\rceil+1$, the time of the queue constraint at slot $j$
  (Definition~\ref{def:good}). On $\Good(<t_j)\cap\bigcap_{I\in\mathcal I_x,\,t(I)\le t_j}\qFr(I)$, with $s_j\in[x,2x)$,
\begin{equation}\label{eq:queueB}
Q_j\le B(x):=\frac{4M_1L^2}{c'}+\frac{16M_1L}{3}+4\sqrt{M_1L(K_d+4)}+8M_1L\Big(\frac L\Lambda+1\Big)+K_d+9.
\end{equation}
\end{enumerate}
In particular $B(x)\le5M_1L^2/c'\le34M_1L^2$. So $B\le L^{O(1)}\qHstd^{1/2}$ off the raised range, and
$B\le X_0^{7.6}L^2$ on it (Theorem~\ref{thm:caplw}). Moreover $B\le H$ at every scale: $\log M_1\le\tfrac12\log H$ since
$H\ge M_1^2$, and $34L^2\le\qHstd^{1/2}=e^{4\ell^2}$, so $B\le(M_1^2\qHstd)^{1/2}\le H$.
\end{theorem}

\begin{proof}[Proof of (i)]
Peel over the value of $\qljit(I)$. For each deterministic $\ell\in\{0,\dots,L_0+5\}$, apply Freedman's inequality in the
form \eqref{eq:freedman}, with deterministic $a(\ell):=\sqrt{2v_\ell\lambda}+\tfrac43M_1\lambda$, $v:=v_\ell$ and
$b:=2M_1$ (by truncation).
\begin{itemize}
\item On $\{\qljit(I)=\ell\}\cap\Good(<t(I))$, the predictable variance satisfies
  $V\le2M_1\qComp_{\mathrm{jit}}(I)\le2M_1(\ell+K_d)=v_\ell$, by R1$'$ (Corollary~\ref{cor:r1}) applied to $I$. So (a) fails
  there only on the event $\{S\ge a(\ell),\,V\le v_\ell\}$, which has probability at most $e^{-\lambda}$.
\item Windows with $\qljit(I)>L_0+5$ are never used in (ii).
\item For (b): take $b=1$. On the good event,
  $\sum_q\Prob(X^{(2)}_q(I)\ge1\mid\mathcal F_{q-})\le\qComp^{(2)}(I)\le s_2(x)|I|\ll1$. This is the source $s_2$ of
  \eqref{eq:s2}, restated in Corollary~\ref{cor:r1}, with \[s_2\le9L\exp\big(-\max(\Lambda,\tfrac12\sqrt{L\log L})\big).\] Since
  $\log(12L\,L_0)\le\log M_1+4\ell+10.2$ and $\log M_1+5\ell+13\ll\tfrac12\sqrt{L\log L}$ at every scale (Theorem~\ref{thm:caplw}),
  $s_2\cdot12L\,L_0\ll1$. Freedman's
  inequality then gives $\#\{q:X^{(2)}_q\ge1\}\le1+\sqrt{2\lambda}+\tfrac23\lambda\le\lambda$, except with probability
  $e^{-\lambda}$.
\item Count: $\#\mathcal I_x\le(x+12L+2)(12L(L_0+1)+3)$, and $e^{-\lambda}=x^{-4}$. \qedhere
\end{itemize}
\end{proof}

\begin{proof}[Proof of (ii)]
\emph{Sandwich.} Let $W\in\mathcal W_x$ have physical interval $[s_{k+1},s_{j+1})$, and let $I\in\mathcal I_x$ be the smallest
window with integer endpoints containing it. Then:
\begin{itemize}
\item $\qCjit(W)\le\sum_q\big(X^{(1)}_q(I)+X^{(2)}_q(I)\big)$, since $\qCjit$ is monotone in the window;
\item $\qljit(I)\le\qljit(W)+4$, since each of the two unit edges holds at most $2$ slots (slot spacing at least $0.99$);
\item $t(I)=b=\lceil s_{j+1}\rceil\le t_j$, and $b\le s_j+11.3L+1\le2x+12L+2$. Indeed a single slot has length at most
  $11.3L$ by \eqref{eq:G1}; slot lengths up to $2$ occur at primes in $(X_0,2X_0)$, and longer ones where $T'_{\mathrm{cont}}<\tfrac12$.
  So $I\in\mathcal I_x$.
\end{itemize}
Hence $\qexc(W)\le\qCjit(I)-\qljit(I)+6$, by the excess identity (Lemma~\ref{lem:excess}).

\emph{Truncation is inactive.} On $\qEcl^{(x)}$ (whose time $8x/L^{11}+H$ is less than $s_j$), Corollary~\ref{cor:increment} applies
to $I$, since $|I|/q\le(12L(L_0+1)+2)/y\le13/(HL^5)\ll\delta_1$. It gives $X^{(1)}_q(I)(U)\le2M_1$ for every value $U$ of the
jitter offset of $q$. So $\tilde X_q\equiv X^{(1)}_q(I)$, and the compensators agree.

\emph{Short windows.} Let $\ell:=\qljit(I)\le\qljit(W)+4\le L_0+4$. By R1$'$,
$\qComp_{\mathrm{jit}}(I)\le\ell(1-c'/L)+K_d$. Together with (a),
\[
\qexc(W)\le-\frac{c'\ell}{L}+\sqrt{16M_1L(\ell+K_d)}+\frac43M_1\lambda+(\text{the $a\ge2$ part})+K_d+6.
\]
The $a\ge2$ part is at most $\lambda\cdot2a_{\max}M_1\le8M_1L(L/\Lambda+1)$, by (b), Corollary~\ref{cor:increment} and
$a_{\max}\le L/\Lambda+1$. Maximising over $\ell\ge0$ gives $\qexc(W)\le B(x)$.

\emph{Full blocks.} If $\qljit(W)=L_0$, then $\ell\ge L_0$ and
\begin{equation}\label{eq:fullblock}
\frac{c'L_0}{L}\ge\sqrt{16M_1L(L_0+4+K_d)}+\frac{16M_1L}{3}+8M_1L\Big(\frac L\Lambda+1\Big)+L^3+K_d+10.
\end{equation}
Indeed, since $M_1\ge c\ge\tfrac14\qHstd^{1/2}\gg L^2$, we have $L_0=48M_1L^3/c'^2$ at every scale. Then
$c'L_0/L=48M_1L^2/c'$ and $\sqrt{16M_1LL_0}=\sqrt{768}\,M_1L^2/c'\le27.72M_1L^2/c'$, so the left side minus the first term
on the right is at least $20.28M_1L^2/c'\ge135M_1L^2$, which exceeds the rest. So $\qexc(W)\le0$ on full blocks.

\emph{Lindley.} Let $k':=\max\{k<j:\ \qljit((k,j])=L_0\}$, if it exists. By \eqref{eq:lindleymax},
\begin{align*}
Q_j&\le\max\Big(Q_{k'}+\qexc((k',j]),\ \max\{\qexc(W):\ W\in\mathcal W_x,\ W=(k,j],\ \qljit(W)<L_0\}\Big)\\
&\le\max(B(s_{k'}),B(x))=B(x).
\end{align*}
This uses: additivity of $\qexc$ over slots; $Q_{k'}\le B(s_{k'})$ on $\Good(<t_j)$, which holds because
$t_{k'}\le s_{k'+1}+2\le s_j$ (a full block has physical length at least $0.99L_0\gg12L$); and that $B$ is non-decreasing
across scales, because $M_1$ and $H$ are non-decreasing in $x$ (Theorem~\ref{thm:caplw}(c)). A window $(k,j]$ with
$\qljit>L_0$ splits at $k'$.
\end{proof}

\begin{remark}[Base scales]\label{rem:basescales}
% v3: §4.3, after the proof of Theorem 4.3
R1$'$ is used on windows $I\subset[x-12L(L_0+1)-2,\,2x+12L+2]$ from $x=X_0$ on. Below $2\Xm$ there are no jittered
composites, so $\qCjit=\qComp_{\mathrm{jit}}=0$ and R1$'$ is trivial there. On $[2\Xm,2X_0)$ it is part of Corollary~\ref{cor:r1}.
\end{remark}

\subsection{The probability skeleton}\label{sec:queue-skeleton}

\begin{definition}[Good event; first failure time]\label{def:good}
% v3: Definition 4.4
Every event $E$ below carries a time $t(E)$, which is a slot, and is $\mathcal F_{t(E)}$-measurable. Put
\begin{gather*}
\Good(<t):=\{Q_j\le B(s_j)\ \text{for all $j$ with $t_j<t$}\}\cap\bigcap_{E:\ t(E)<t}E,\\
t_j:=\lceil s_{j+1}\rceil+1,\qquad\tau:=\min\{t:\ \Good(\le t)\text{ fails}\},
\end{gather*}
with $B(s):=B(x)$ for $s\in[x,2x)$. Then $\tau$ is an $(\mathcal F_t)$-stopping time. The events are listed per dyadic scale
$u=2^mX_0$, with parameters frozen at $u$, in Table~\ref{tab:events}.
\end{definition}

\begin{table}[ht]
\centering\small
\renewcommand{\arraystretch}{1.25}
\begin{tabular}{@{}p{0.17\textwidth}p{0.33\textwidth}p{0.25\textwidth}p{0.16\textwidth}@{}}
\hline
event & windows at scale $u$ & $\Prob(E^c\cap\Good(<t(E)))$ & $t(E)$\\
\hline
E1: $\qEcl^{(u)}$ (Proposition~\ref{prop:ecl}) & grid windows, left end in $[X_0/2,\,4u/L^{11}]$ & $\le L_u^{-2}$ & $8u/L^{11}+H$\\
E2: $\qFr(I)$ (Theorem~\ref{thm:queue}) & grid windows $I\in\mathcal I_u$ (integer endpoints, right end $\le2u+12L+2$, length
  $\le12L(L_0+1)+2$), peeled over $\qljit(I)$ & $\le u^{-2}$ & $b=\max I$\\
E3: $G^{\mathrm{NV}}_u$ (Lemma~\ref{lem:noise}), including the bound \eqref{eq:F2} & windows with integer endpoints in $[u,2u]$,
  length in $[H-1,\,1.03H+1]$, two-sided, $\lambda'_u:=3\log u+2\log\log u+2+\log H(u)$ & $\le1.2u^{-2}L_u^{-2}$ for E3 and
  E3$'$ together (on $\qEcl(u)$; sure variance bound $V_I\le2M_1\qComp(I)$) & each window when it closes, $\le2u$\\
E3$'$: $G^{\mathrm{sw}}_u$ (Lemma~\ref{lem:noise}), including \eqref{eq:F2} & windows with integer endpoints in $[u,2u]$, length
  $\le1.03H+1$, $\lambda'_u$ & (with E3) & each window when it closes, $\le2u$\\
E4: $\qLand(J)$ (Lemma~\ref{lem:landing}) & landing windows $J\in\mathcal J_u$ (half-integer grid in $[u,4u+1]$, length
  $\le6H(u)$) & $\le146H(u)u^{-3}$ & $\max J+H(\max J)$\\
\hline
\end{tabular}
\caption{The events of Definition~\ref{def:good} at the dyadic scale $u=2^mX_0$.}\label{tab:events}
\end{table}

There is no separate event for small latest primes: since $\Lambda\ge10^8>L_*$, these are absorbed (Section~\ref{sec:queue-abs}), and
the small-$q$ case of S1 needs $L<L_*$ (Proposition~\ref{prop:S1S2}).

\begin{lemma}[First failure; probability budget]\label{lem:budget}
% v3: Lemma 4.5
\[
\Prob(\tau<\infty)\le\sum_{m\ge0}\ \sum_{E\text{ at scale }2^mX_0}\Prob(E^c\cap\Good(<t(E)))\le\frac{1.5}{\Lambda}<\frac12.
\]
On $\{\tau=\infty\}$, $Q_j\le B(s_j)$ for all $j$, every event of Table~\ref{tab:events} holds, and \eqref{eq:La} holds pathwise:
$|D(x)|\le B(x)+\tfrac{H(x)}2(1+\kappa_+/\log x)+3\le x^{\eta}$ for all $x\ge X_0+2H$.
% Approved deviation from v3 (CORRECTIONS-vs-v3.md, #15): v3 has H(x)/2+2.
\end{lemma}

\begin{proof}
\emph{($*$) Times of the events.} Each bound in Table~\ref{tab:events} holds on $\Good(<t(E))$. No conditioning is used: Freedman's
inequality is applied to deterministic windows, with deterministic $(a,v)$ after peeling over $\qljit$ (E2; likewise E3 and
E3$'$ in Lemma~\ref{lem:noise}). Every consumer of an event $E$ runs at a time at least $t(E)$:
\begin{itemize}
\item E1 is used by Theorem~\ref{thm:queue}, by the noise bounds of Lemma~\ref{lem:noise} and by Lemma~\ref{lem:lemmaP}, at scales $\ge u$;
\item E2 is used by the queue constraint at slot $j$, whose time $t_j$ is at least $t(I)$;
\item E3 is used by $\nu$ on $[u,2u]$, read at times $\ge2u$;
\item E3$'$ is used by Proposition~\ref{prop:pointwise} (the pointwise margin);
\item E4 is used by the local counts, read at times at least $4$ times their height.
\end{itemize}

\emph{(D) The queue constraint at slot $j$ carries the time $t_j:=\lceil s_{j+1}\rceil+1$.}
\begin{itemize}
\item $t_j$ is predictable, because $T$ is adapted (Lemma~\ref{lem:adapted}). The outer window $I\supset W$ of the proof of
  Theorem~\ref{thm:queue}(ii) ends at $b=\lceil s_{j+1}\rceil<t_j$, and R1$'$ on $I$ holds on $\Good(<\max I)$.
\item $\Good(<t_j)\cap\bigcap_{t(E)=t_j}E\subset\{Q_j\le B(s_j)\}$ by Theorem~\ref{thm:queue}(ii), and the inputs of that
  theorem have times at most $t_j$.
\item $t_j\le s_j+11.3L+2$. Every consumer of a queue constraint reads it with margin at least $H\gg12L$:
  Corollary~\ref{cor:localcount} at height $z$ reads slots $\le z+H$ and runs after $z+7H$; the Lindley step reads $Q_{k'}$ with
  $t_{k'}\le s_j$; and \eqref{eq:La}, Section~\ref{sec:zeros} and Section~\ref{sec:noise} read queue bounds at heights $\le x/2$ or with
  margin at least $H/2$.
\item No event with time at most $t_j$ reads $Q_j$.
\end{itemize}

\emph{First failure.} On $\{\tau=t\}$, $\Good(<t)$ holds and $\Good(\le t)$ fails. If some $E$ with $t(E)=t$ fails, then
$E^c\cap\Good(<t(E))$ occurred. Otherwise the queue constraint with time $t$ fails, which (D) excludes.

\emph{Summation.} Summing Table~\ref{tab:events} over $m$, E1 contributes
$\sum_mL_u^{-2}\le\sum_{m\ge0}(\Lambda+m\log2)^{-2}\le1.45/\Lambda$. Everything else contributes
$\le\sum_uO(u^{-2})\ll1/X_0$.

The target $T$ is adapted (Section~\ref{sec:construction}), and every bound above holds for any adapted $T$. The union runs over
dyadic scales, with finitely many windows per scale, never over all $x$; note that $\sum_x1/x$ diverges.
\end{proof}

\subsection{The cluster event and its reduction}\label{sec:queue-ecl}
% v3: §4.4

Fix a dyadic height $z\ge X_0/2$; this covers the band $[\Xm,X_0)$ of centred placement. Put $L:=\log z$, and let $J$ be a grid
window at height $z$, so $|J|=2\delta_1=4L^{-6}<1$. Since $|J|<1$, $J$ contains at most one element with no jittered prime.
Write $S(J)$ for the set of elements of $\NN$ in $J$ that have a jittered prime. For $n=M\cdot d$, with $M$ the jittered part
and $d$ the rigid ($X_0$-smooth) part, the \emph{exponent vector} $f_n\in\Z_{\ge0}^{\mathcal P_{\mathrm{jit}}}$ of $n$ is
that of $M$. \emph{Rank}, \emph{span} and \emph{degree} $\qdeg f:=\|f\|_1$ refer to these vectors. Every degree is at most
$D:=\log2z/\log\Xm$, which equals $\log2z/\Lambda$ up to a relative error $O(H(X_0)/X_0)$.
\begin{itemize}
\item \emph{(V1) One element per vector.} Two elements of $J$ with the same exponent vector $f$ have the same jittered part
  $M_f\ge X_0$. They differ by $M_f|d-d'|\ge X_0>|J|$ unless $d=d'$. So distinct elements of $S(J)$ have distinct exponent
  vectors.
\item \emph{Capacity.} For $S'\subset S(J)$ of rank $r$, put $\qCap(S'):=|S'|-r$.
\end{itemize}

\begin{proposition}[The cluster event]\label{prop:ecl}
% v3: Proposition 4.6
Let $\qEcl(z)$ be the event that every grid window $J$ at height $z$ satisfies $\#(\NN\cap J)\le M_1(z)$. Then, with
$t_z:=2z+H(2z)$,
\[
\Prob(\qEcl(z)^c\cap\Good(<t_z))\le L^{-2}.
\]
The event $\qEcl^{(u)}:=\bigcap_{X_0/2\le z\le4u/L_u^{11}}\qEcl(z)$ (over dyadic $z$) is the event E1 of Definition~\ref{def:good}. Its
new heights at scale $u$ form one dyadic block, so its failure probability on the good event is at most
$L_u^{-2}(1+o(1))$.
\end{proposition}

\begin{proof}
Suppose $\#(\NN\cap J)\ge M_1+1$, so $|S(J)|\ge M_1=r_0+\varphi(r_0)+1$. Let $r:=\operatorname{rank}S(J)$.
\begin{itemize}
\item \emph{Low heights, $2z<(\Xm)^2$}, in particular every $z\in[X_0/2,X_0)$. Every element of $S(J)$ is $q\cdot d$ with
  exactly one jittered prime $q$, because two jittered primes give a product at least $(\Xm)^2>2z$. (Heights $z\ge X_0$ are in
  any case covered by the general argument below.) So the exponent vectors are unit vectors, distinct by (V1). Hence
  $\operatorname{rank}S(J)=|S(J)|$ and $\qCap=0$, so the case $r<r_0$ below is vacuous. The case $r\ge r_0$ is the direct
  (F1)--(F2) count $E\le(1.01|J|)^{r_0}/r_0!$, to which Theorem~\ref{thm:R} applies; Theorem~\ref{thm:R} is stated for $z\ge X_0/2$.
\item \emph{Case $r\ge r_0$.} $S(J)$ contains a rank-full $r_0$-subset. By Theorem~\ref{thm:R} with $k=r_0$ and
  $\delta=|J|=4L^{-6}$, the expected number of rank-full $r_0$-subsets of $\NN\cap J$ on $\Good(<t_z)$ is at most
  $(\theta A_R\delta)^{r_0}e^{(3H(X_0)+2)r_0/X_0}(1+r_0/X_0)^{1.02r_0}e^{2.2r_0/\Lambda}$; the factor $e^{(3H(X_0)+2)r_0/X_0}$
  is absorbed in the per-element factor $1.01$. There are at most $2z/\delta$ windows. The total is at most $L^{-4}$ for
  $L\ge465$, by Corollary~\ref{cor:R-Ecl}.
  % Approved deviation from v3 (CORRECTIONS-vs-v3.md, #1): v3 has e^{2r_0/X_0} in place of e^{(3H(X_0)+2)r_0/X_0}.
\item \emph{Case $r<r_0$.} Then $|S(J)|>r_0+\varphi(r_0)\ge r+\varphi(r)$. So $J$ is a window as in
  Theorem~\ref{thm:caplw}(a): with $V$ a heavy subspace of $S(J)$ of least dimension, Lemma~\ref{lem:heavy} gives
  $\qCap(S(J)\cap V)\ge c$. By Theorem~\ref{thm:caplw}(a), the probability that some window at height $z$ is of this kind
  is at most $L^{-3}$.
\end{itemize}
Summing the two cases gives the claim.
\end{proof}

\begin{corollary}[Per-prime increment bound]\label{cor:increment}
% v3: Corollary 4.7
Let $u\ge X_0$ be a height. On $\qEcl^{(u)}$ the following holds for every window $I\subset[u,2u+12L_u+2]$ (slot window or
physical window) with $|I|\le\max(H(u),12L_u(L_0(u)+1)+2)$, every non-absorbed latest prime $q\ge y(u)$, every $a\ge1$ and
\emph{every} value of the jitter offset of $q$:
\[
X^{(a)}_q(I):=\#\{n\in I\setminus\Abs:\ q^a\,\|\,n,\ q\text{ latest}\}\le2M_1(u).
\]
\end{corollary}

\begin{proof}
The cofactors $m=n/q^a$ lie in $I/q^a$, which has length at most
\[
|I|/y\le\max\big(1/(L_0L^6),\,13/(HL^5)\big)\le\delta_1(u/q),
\]
since $y=H\cdot L_0\cdot L^6$. They are $q$-smooth, so they are decided no later than $q$.
\begin{itemize}
\item If $I/q^a$ meets $[\Xm,\infty)$, it lies at height at least $X_0/2$. It is covered by two grid windows of $\qEcl$ at that
  height, each holding at most $M_1(u/q^a)\le M_1(u)$ elements, since $M_1$ is non-decreasing.
\item If $I/q^a$ lies below $\Xm$, the cofactors have no jittered prime, so they are rigid integers. Since $I/q^a$ has length
  less than $1$, it holds at most one of them. \qedhere
\end{itemize}
\end{proof}

\begin{remark}\label{rem:increment-uses}
% v3: §4.4, after Corollary 4.7
Corollary~\ref{cor:increment} is all that Section~\ref{sec:noise} (Lemma~\ref{lem:noise}, Proposition~\ref{prop:pointwise}) and Theorem~\ref{thm:queue} take from
$\qEcl$. Shorter windows only make it easier. It covers every window of length at most $L^{-6}$ at any height $\ge X_0$,
including the cofactor windows at height $u/q^2$ of the $q^2$-composites in Section~\ref{sec:noise}.
\end{remark}

\begin{lemma}[Heavy subspace]\label{lem:heavy}
% v3: Lemma 4.8 (Participation+)
Let $S\subset\Z^m_{\ge0}\setminus\{0\}$ be finite, of rank less than $r_0$, with $|S|\ge r_0+\varphi(r_0)$. Call a subspace
$W$ spanned by elements of $S$ \emph{heavy} if $\#(S\cap W)\ge\varphi(\dim W)$. Let $V$ be a heavy subspace of least
dimension $r$. Then:
\begin{enumerate}[label=\textup{(\roman*)}]
\item $\qCap(S\cap V)\ge\varphi(r)-r\ge c$;
\item every proper subspace $W'\subsetneq V$ spanned by elements of $S$ satisfies $\#(S\cap V\setminus W')\ge Q+r+1$;
\item if the elements of $S\cap V$ of degree at most $k$ number at most $Q$, then the elements of $S\cap V$ of degree at
  least $k+1$ span $V$.
\end{enumerate}
\end{lemma}

\begin{proof}
$\operatorname{span}S$ is heavy, since $\varphi$ is increasing and $|S|\ge\varphi(r_0)\ge\varphi(\operatorname{rank}S)$. So
$V$ exists, and (i) follows from $\varphi(r)\ge c+r$.

(ii) $W'$ is not heavy, by minimality of $V$, so $\#(S\cap W')\le\varphi(\dim W')-1\le\varphi(r-1)-1$; and
$\varphi(r)-\varphi(r-1)=Q+r$.

(iii) Otherwise the elements of high degree lie in a proper element-spanned subspace $W'$, so $S\cap V\setminus W'$
consists of elements of low degree only, of which there are at most $Q<Q+r+1$. This contradicts (ii).
\end{proof}

\subsection{Capacity forces thin covers}\label{sec:queue-cover}
% v3: §4.5 (deterministic combinatorics)

Throughout this subsection $S\subset\Z^m_{\ge0}\setminus\{0\}$ is finite, of rank $r$, and every vector in $S$ has
$\ell^1$-norm at most $D$. Put $P(S):=\bigcup_{v\in S}\qsupp v$. A \emph{cover} is a subset $C\subset S$ with
$\bigcup_{y\in C}\qsupp y=P(S)$, and $\kappa(S)$ is the least size of a cover. For a basis $B\subset S$ and $y\in S\setminus B$
(an \emph{extra}), write $e_y=\sum_{i\in B}\alpha_ie_i$, $I_+:=\{i:\alpha_i>0\}$, $I_-:=\{i:\alpha_i<0\}$ and
$\qsupp_B(y):=I_+\cup I_-$. No affine-hyperplane assumption is used in this subsection, except in the last statement of
Theorem~\ref{thm:tcl}.

\begin{lemma}[Covering lemma; minimal covers]\label{lem:covering}
% v3: Lemma 4.9
\begin{enumerate}[label=\textup{(\alph*)}]
\item Every basis $B\subset S$ is a cover.
\item If $S$ has an extra $y$ that is not proportional to a basis vector, then $\kappa(S)\le r-1$.
\item A minimal cover $C$ (one with no proper subset that covers) is linearly independent.
\end{enumerate}
\end{lemma}

\begin{proof}
(a) Every $e_y$ is a combination of the $e_i$, $i\in B$, and nonnegativity gives
$\qsupp e_y\subset\bigcup_{\alpha_i\ne0}\qsupp e_i$.

(b) Rearrange as $e_y+\sum_{l\in I_-}|\alpha_l|e_l=\sum_{i\in I_+}\alpha_ie_i$. All coefficients are positive and all vectors
are nonnegative, so
\begin{equation}\label{eq:supportid}
\qsupp(y)\cup\bigcup_{l\in I_-}\qsupp(e_l)=\bigcup_{i\in I_+}\qsupp(e_i).
\end{equation}
If $I_-\ne\emptyset$, then $B\setminus I_-$ covers. If $I_-=\emptyset$, then $\{y\}\cup(B\setminus I_+)$ covers, and it has
size at most $r-1$ when $|I_+|\ge2$. The remaining case, $|I_+|=1$ with $I_-=\emptyset$, is the proportional case.

(c) Each $c\in C$ has a prime $p_c$ that lies in no other element of $C$, for otherwise $C\setminus\{c\}$ covers. The
$C\times\{p_c\}$ submatrix is diagonal with positive diagonal.
\end{proof}

\begin{lemma}[Coefficient support]\label{lem:coefsupport}
% v3: Lemma 4.10
If $\kappa(S)\ge r-j$, then for every basis $B$ and every extra $y$ we have $|I_-|\le j$ and $|I_+|\le j+1$. Hence
$|\qsupp_B(y)|\le2j+1$.
\end{lemma}

\begin{proof}
If $|I_-|\ge1$, then $B\setminus I_-$ covers by \eqref{eq:supportid}, so $r-|I_-|\ge\kappa\ge r-j$. In all cases
$\{y\}\cup\{e_l:l\in I_-\}\cup(B\setminus(I_+\cup I_-))$ covers by \eqref{eq:supportid}. It has size $r+1-|I_+|$, so
$|I_+|\le j+1$.
\end{proof}

\begin{theorem}[Capacity forces thin covers]\label{thm:binomial}
% v3: Theorem 4.11
For every $j\ge0$,
\[
\kappa(S)\ge r-j\ \Longrightarrow\ \qCap(S):=|S|-r\le N(r,D,j):=\sum_{s=1}^{2j+1}\binom rs\binom{D+s}s.
\]
Equivalently, $\qCap(S)>N(r,D,j)$ implies $\kappa(S)\le r-1-j$. Also $N(r,D,j)\le e\,(r(D+1))^{2j+1}$.
\end{theorem}

\begin{proof}
Fix a basis $B$. By Lemma~\ref{lem:coefsupport}, each extra lies in $\operatorname{span}\{e_i:i\in I\}$ for some $I\subset B$ with
$|I|=s\le2j+1$. Choose $s$ coordinates on which this span projects injectively. The extras with support $I$ then map
injectively into $\{v\in\Z^s_{\ge0}:\|v\|_1\le D\}$, which has exactly $\binom{D+s}s$ points. Sum over $I$.
\end{proof}

\begin{theorem}[Thin-Cover Lemma]\label{thm:tcl}
% v3: Theorem 4.12
If $\kappa(S)\ge r-1$, then $\qCap(S)\le4r(D+1)^3$. If $S$ lies on an affine hyperplane $w^{\mathsf T}e=1$ not through $0$,
the bound improves to $\qCap(S)\le3r(8D+4)$.
\end{theorem}

This is the case $j=1$ of Theorem~\ref{thm:binomial}, with $O(r)$ supports instead of $\binom r3$. It is recorded for
completeness; Section~\ref{sec:queue-capl} uses Theorem~\ref{thm:binomial}.

\begin{proof}
\emph{Supports have size at most $3$.} This is Lemma~\ref{lem:coefsupport} with $j=1$. A support of size $1$ means that $y$ is
proportional to a basis vector; there are at most $r$ such supports, each holding at most $D+1$ points.

\emph{Exchange rule.} Write $I(y):=\qsupp_B(y)$. Let $I(y)\cap I(y')=\{h\}$ for extras $y,y'$. Put $B':=B-e_h+y$ and
substitute for $e_h$ in $y'$. Nothing cancels, so $|\qsupp_{B'}(y')|=|I(y)|+|I(y')|-1\le3$, and hence
$|I(y)|=|I(y')|=2$. In particular a path of three $2$-supports $\{1,2\},\{2,3\},\{3,4\}$ is impossible: after exchanging
$e_2$ for $y_{12}$ it produces a $3$-support meeting a $2$-support in exactly one point.

\emph{Components.} In the hypergraph of supports of sizes $2$ and $3$, each component on $v_c$ vertices of $B$ is one of:
a star or a triangle of pairs; a book $\{i,j,l_t\}$, possibly together with the pair $\{i,j\}$; or a family inside a $4$-set.
Each has at most $3v_c$ supports, so there are at most $3r$ in all.

\emph{Points per support.} A support of size $s$ contributes at most $(D+1)^s$ points, as in Theorem~\ref{thm:binomial}. So
$\qCap\le r(D+1)+3r(D+1)^3\le4r(D+1)^3$.

\emph{Hyperplane case.} By Lemma~\ref{lem:coefsupport} with $j=1$, $S$ intersected with a plane is weakly convex: no point of $S$
lies strictly inside a triangle of points of $S$. Projecting to $[0,D]^2$ gives at most $4(2D+1)$ points per plane.
\end{proof}

\subsection{High-capacity configurations are rare at every scale}\label{sec:queue-capl}
% v3: §4.6 (CapL-w)

\paragraph{Setting.} Fix a dyadic height $z\ge X_0$, and put $L:=\log z$, $\ell:=\log L$ and $D:=\log2z/\Lambda$. Let $J$ range
over the $2z/\delta$ grid windows at height $z$, where $\delta:=2\delta_1=|J|$. A \emph{configuration} is a rank-full set
$B'\subset S(J)$ of rank $r<r_0$, together with a minimum cover $C\subset B'$ of $P(B')$, so $|C|=\kappa$. Write $R:=B'\setminus
C$ for the \emph{rows outside the cover}. For $y\in B'$ write $n_y=M_yd_y$, with $M_y$ jittered and $d_y$ rigid. The
\emph{type} $\tau$ of a configuration is its exponent pattern, that is, the matrix of exponents of its distinct jittered
primes, up to relabelling by value. Let $\pi:=|P(C)|\le\kappa D$ be the number of jittered primes, and put
$\lambda_D:=\sum_{X_0\le q\le2z}1/q$; on the good event, $\lambda_D\le\log D+1$ by (F4).

\begin{proposition}[Cover Bound]\label{prop:cover}
% v3: Proposition 4.13
For every type $\tau$,
\begin{equation}\label{eq:coverbound}
\E\big[\#\{\text{configurations of type $\tau$ in $J$}\};\ \Good(<t_z)\big]\le(\theta\delta)^r\sum\Big(\prod_{y\in C}d_y^{-1}
\cdot X_0^{-\mathrm{exc}}\cdot\prod_{f\in F}q_f^{-1}\cdot\prod_{y\in R}n_y^{-1}\Big),
\end{equation}
where:
\begin{itemize}
\item the sum runs over the realised jittered primes $F$ that are not used as pivots, over the tails of the rows of $C$, and,
  for each $y\in R$, over its monomial $M_y$ in the labels of $P(C)$ and its tail $d_y$, subject to $n_y\in J$;
\item $\mathrm{exc}:=\sum_p(\#\{\text{occurrences of $p$ in rows of $C$}\}-1)\ge0$.
\end{itemize}
\end{proposition}

\begin{proof}
\emph{Exact identity.} By the Volume Lemma (Lemma~\ref{lem:volume}), the pivot-integral of $\prod_{y\in B'}\mathbf 1(n_y\in J)$
equals $\int_{J^r}\prod_{\mathrm{piv}}q\big/\big(|\det E_P|\prod_yn_y\big)\,dn$.
\begin{itemize}
\item Order the pivots latest first, and stage the rows by Lemma~\ref{lem:SC} (nested lexicographically-first bases with the
  rows of $C$ listed first).
\item The Staged Volume Lemma (Theorem~\ref{thm:svl}) then integrates the pivots stage by stage. Every late free block is handled by
  a block bound (the realised block sum, bounded as in Lemma~\ref{lem:fac}\,(S3)) against a weight that is harmonic in
  some row of $C$ of its stage. This is the condition that Lemma~\ref{lem:SC}
  guarantees.
\item $\prod_{\mathrm{piv}}q/\prod_yn_y=\prod_yc_y^{-1}\cdot\prod_{\mathrm{piv}}q^{1-s_p}$, where $c_y:=n_y/(\text{pivot part
  of }n_y)$ and $s_p$ is the column sum of $p$.
\end{itemize}

\emph{Allocation.} Every prime has an occurrence in some row of $C$, since $C$ covers $P(B')$. So:
\begin{itemize}
\item each pivot $p$ cancels the factor $q_p^{+1}$ of the Jacobian against one occurrence chosen in a row of $C$;
\item each free prime $f$ keeps one occurrence in a row of $C$ for its harmonic sure sum (Lemma~\ref{lem:fac}, (S1$'$) and (S3));
\item all occurrences in rows of $R$ stay with their row;
\item the remaining occurrences in rows of $C$ are the excess, each at most $X_0^{-1}$.
\end{itemize}
The resulting weights are as follows. A row of $C$ gets at most $d_y^{-1}$. A row of $R$ gets exactly
$c_y^{-1}\cdot\prod_{p\in M_y}p^{-E}=1/n_y$; its tail cancels, so it is summed with counting weight, which is a deterministic
rigid count. Free primes are harmonic, and the excess contributes $X_0^{-\mathrm{exc}}$.

\emph{Bookkeeping.} The factor $\theta=1+1/\Lambda$ per element comes from (F1)--(F2) and the slot sums, and $|\det E_P|\ge1$.
\end{proof}

\begin{remark}\label{rem:jacobianrows}
% v3: §4.6, Remark after Proposition 4.13
Every row of the basis $B'$ is a Jacobian row. Only extras outside $B'$ would be mere indicators, and none occur
(Lemma~\ref{lem:reduction}(i)). Being a Jacobian row is what supplies the factor $1/n_y\le1/z$ of a row of $R$, and it does not
cost the rigid count. The per-row weights below never credit $X_0^{-\mathrm{exc}}$ beyond the pattern sum over the rows of
$C$ in Lemma~\ref{lem:reduction}(ii).
\end{remark}

\begin{lemma}[Reduction to configurations; type union]\label{lem:reduction}
% v3: Lemma 4.14
\begin{enumerate}[label=\textup{(\roman*)}]
\item (Extras drop out.) Let $S':=S(J)\cap V$ with $V$ as in Lemma~\ref{lem:heavy}, of rank $r$ and with $\qCap(S')\ge c$. Then
  $\kappa(S')\le r-1-j(c)$, where $j(c):=\max\{j:\ e(r(D+1))^{2j+1}<c\}$ (Theorem~\ref{thm:binomial}). A minimum cover $C$ of $S'$
  is independent (Lemma~\ref{lem:covering}(c)), so it extends to a basis $B'\subset S'$ with $P(B')=P(S')$
  (Lemma~\ref{lem:covering}(a)). So $|R|=r-\kappa\ge1+j(c)$. If moreover the elements of $S'$ of degree at most $k$ number at most
  $Q$, the completion can be chosen with every row of $R$ of degree at least $k+1$ (Lemma~\ref{lem:heavy}(iii) and Steinitz
  exchange).
\item (Type union.) Sum \eqref{eq:coverbound} over types and multiply by the $2z/\delta$ windows. For $L\ge60$, the result is,
  per height, at most
  \begin{equation}\label{eq:typeunion}
  13L^{5.1}\cdot z\cdot\prod_{y\in R}w_y,\qquad w_y:=z^{-1}\sum_{M_y}\Psi(2z/M_y,X_0),
  \end{equation}
  where $M_y$ runs over the monomials of the degree class of the row in the $\pi$ labels.
\item (Consequence.) The probability that some window at height $z$ contains a configuration with $|R|\ge1+j$ and all
  $w_y\le w$ is at most $r_0\cdot13L^{5.1}z\cdot w^{1+j}$. It is at most $L^{-3}$ as soon as
  \begin{equation}\label{eq:criterion}
  (1+j)\log(1/w)\ge L+10\ell+10.
  \end{equation}
\end{enumerate}
\end{lemma}

\begin{proof}
\emph{Proof of (ii).} Per row, the type and sum factors are:
\begin{itemize}
\item the free primes of a row of $C$, sorted by value: $\sum_k\lambda_D^k/k!\le e^{\lambda_D}\le eD$;
\item the pivot positions: at most $D$ per row;
\item the excess patterns of the rows of $C$: the sum over patterns of $X_0^{-\mathrm{exc}}$ is at most
  $\prod_{p\in P(C)}(1-1/p)^{-\kappa}\le L^{2.1}$ per row of $C$, using that the realised primes $p\ge X_0$ are distinct and
  $\pi\le L^2$;
\item the tail of a row of $C$: $\sum_d1/d\le e^{\gamma}\Lambda(1+1/\Lambda^2)\le2L$.
\end{itemize}
The minimum cover used is a fixed choice determined by the type, so its choice costs nothing. So each row costs at most
$C':=6L^{5.1}$, and $(2z/\delta)(\theta\delta C')^r\le2z\cdot\theta C'\cdot(\theta\delta C')^{r-1}\le13L^{5.1}z$, since
$\theta\delta C'=24\theta L^{-0.9}\le1$ for $L\ge60$. The monomials and tails of the rows of $R$ are inside $w_y$.

\emph{Proof of (iii).} Sum (ii) over the fewer than $r_0$ possible ranks. Since $w\le1$, we have
$\prod_{y\in R}w_y\le w^{|R|}\le w^{1+j}$. As $z=e^L$, $r_0\le L=e^{\ell}$ and $\log13\le2.6$, the bound
$r_0\cdot13L^{5.1}z\,w^{1+j}$ is at most $L^{-3}$ once $(1+j)\log(1/w)\ge L+9.1\ell+2.6$, which \eqref{eq:criterion} implies.
\end{proof}

\begin{lemma}[Row weights]\label{lem:rowweights}
% v3: Lemma 4.15
Let $y\in R$ have degree $k\ge1$, so that $M_y\ge X_0^k$ and $d_y\le2z/X_0^k$.
\begin{enumerate}[label=\textup{(\alph*)}]
\item $w_y\le\sum_{k'\ge k}\binom{\pi+k'-1}{k'}\Psi(2z/X_0^{k'},X_0)/z$.
\item (High degree.) If $e(\pi+k)/(kX_0)\le1/3$, then $w_y\le3\,(e(\pi+k)/k)^kX_0^{-k}$; that is,
  $\log(1/w_y)\ge k(\Lambda-1-\log((\pi+k)/k))-\log3$.
\item (Smooth tails.) Put $u:=\log Y/\Lambda$. For $\max(10^4,\log\Lambda)\le u\le X_0^{1/2}$ we have
  $\Psi(Y,X_0)\le Y\exp(-0.9u\log u)$. For $u>X_0^{1/2}$ we have
  \[\Psi(Y,X_0)\le\exp\big(1.26X_0\log(2\log Y)/\Lambda\big).\]
\item (Low degree, smooth tails.) If $10^4\le D-1\le X_0^{1/2}$ and $\pi X_0^{-1}(eD)^{0.9}\le1/3$, then
  $\log(1/w_y)\ge g:=\Lambda-\log(3\pi)+0.9(D-1)\log(D-1)$.
\item (Joint Rankin bound; any $L$.) For every $\eta\in[\tfrac14,1)$,
  \[
  \log(1/w_y)\ge\eta L-\log2-\log\Lambda-2-\frac{1.3X_0^{\eta}}{\eta\Lambda}-\frac{1.36(X_0+\pi)^{\eta}}{0.99\,\eta}.
  \]
\end{enumerate}
\end{lemma}

\begin{proof}
(a) Count the pairs $(M_y,d_y)$.

(b) Use $\Psi(Y,X_0)\le Y$. The terms $\binom{\pi+k'-1}{k'}\cdot2X_0^{-k'}$ have consecutive ratio
$(\pi+k')/((k'+1)X_0)\le1/(3e)$, by the hypothesis. So the sum over $k'\ge k$ is at most
$2.3\binom{\pi+k-1}{k}X_0^{-k}\le3(e(\pi+k)/k)^kX_0^{-k}$.

(c) We use Rankin's bound $\Psi(Y,y)\le Y^{\sigma}\prod_{p\le y}(1-p^{-\sigma})^{-1}$ with $y=X_0$ and $\sigma=1-\eta$, where
$\eta:=\log(u\log u)/\Lambda\in(0,1)$.
\begin{itemize}
\item $Y^{\sigma}=Y\exp(-u\log u-u\log\log u)$.
\item By partial summation with $\pi(t)\le1.26t/\log t$ and $\sum_{p\le y}1/p\le\log\log y+0.27$ for $y\ge2140$
  \cite[Thm.~5, (3.18)]{RosserSchoenfeld1962},
  \[
  \log\prod_{p\le X_0}(1-p^{-\sigma})^{-1}\le\log\Lambda+0.27+1.26\,\frac{u\log u}{\log(u\log u)}\Big(1+\frac2{\log(u\log u)}\Big)+1
  \le\log\Lambda+1.3u+2.
  \]
\item So $\log(\Psi/Y)\le-u\log u-u\log\log u+1.3u+\log\Lambda+2\le-0.9u\log u$ for $u\ge\max(10^4,\log\Lambda)$. This needs
  $0.1u\log u+u\log\log u-1.3u\ge\log\Lambda+2$, which holds there. (At $u=10^4$ alone it would need
  $\log\Lambda\le1.84\cdot10^4$; this is why the lower limit is $\max(10^4,\log\Lambda)$.)
\item For $u>X_0^{1/2}$: an $X_0$-smooth $d\le Y$ has at most $1+\log Y/\log2$ choices of exponent for each of the at most
  $1.26X_0/\Lambda$ primes $p<X_0$.
\end{itemize}

(d) Put $u_1:=\max(10^4,\log\Lambda)$. The $k=1$ term of (a), with (c) at $u=D-1$, is at most
$2\pi\exp(-0.9(D-1)\log(D-1))/X_0$.
\begin{itemize}
\item For $D-k'\ge u_1$, the ratio of consecutive terms is at most $\pi X_0^{-1}(eD)^{0.9}\le1/3$, by hypothesis.
\item For $D-k'<u_1$, bound the term by $\Psi\le Y$ instead. The first such term is at most
  $\pi X_0^{-1}e^{0.9(u_1+1)\log(u_1+1)}$, which is far less than $\tfrac13$ times the last term of type (c), and the
  following ratios are at most $(\pi+k')/((k'+1)X_0)\le\tfrac13$.
\item So $g$ is unchanged. Where (d) is used (Steps 3--4 of the proof of Theorem~\ref{thm:caplw}), $D-1\ge\Lambda e^{s_H+1}\gg u_1$.
\end{itemize}

(e) We have
\[
w_y\le z^{-1}\sum_{(M,d)}\Big(\frac{2z}{Md}\Big)^{1-\eta}\le2z^{-\eta}\prod_{p\in P(C)}(1-p^{\eta-1})^{-1}
\prod_{p<X_0}(1-p^{\eta-1})^{-1}.
\]
\begin{itemize}
\item The rigid product is at most $\exp(\log\Lambda+2+1.3X_0^{\eta}/(\eta\Lambda))$, by partial summation as in (c).
\item The jittered primes of $P(C)$ are $\pi$ distinct primes in distinct slots. The $i$ smallest occupy $i$ distinct
  slots, the largest of which is at least $X_0+0.99(i-1)$. Since the jitter offset satisfies $|q-s|\le H(s)/2\le s/3$
  (see \eqref{eq:C2}), the $i$-th smallest prime is at least $\tfrac23(X_0+0.99(i-1))$. So
  \[
  \sum_{p\in P(C)}p^{\eta-1}\le(3/2)^{1-\eta}\sum_{i=1}^{\pi}(X_0+0.99(i-1))^{\eta-1}\le\frac{1.36(X_0+\pi)^{\eta}}{0.99\,\eta},
  \]
  using $(3/2)^{1-\eta}\le1.36$ for $\eta\ge\tfrac14$. \qedhere
\end{itemize}
\end{proof}

\begin{theorem}[High-capacity configurations at every scale; the capacity parameter]\label{thm:caplw}
% v3: Theorem 4.17 (CapL-w at every scale, uniformly in Lambda >= 10^8)
Let $\Lambda=\log X_0\ge10^8$ be arbitrary (Section~\ref{sec:params}), and let $\theta\in[\tfrac14,\theta_1]$ with
$\theta_1=\tfrac12-10^{-4}$, as in Section~\ref{sec:setup}, so that $\theta_b=\theta/2\ge\tfrac18$. For a dyadic height $z\ge X_0$
put $\ell:=\log\log z$ and $s:=\ell-2\log\Lambda$, so that $L=\Lambda^2e^s$. Define
\begin{align*}
A(\Lambda,s)&:=\log\big(2\Lambda^3e^{2s}/(2\log\Lambda+s)\big)=\log\big(2L^2/(\ell\Lambda)\big),&
G(\Lambda,s)&:=\frac{2\sqrt{A(A+1)}\,e^{s/2}}{\sqrt{\Lambda-A-1}},\\
\Lambda_s&:=\max(10^8,e^{-s}),& \gamma_*(s)&:=1.3\,G(\Lambda_s,s).
\end{align*}
Let $s_H=5.9745\ldots$ be the root of $\gamma_*(s_H)=\tfrac13-10^{-5}$. Put $\gamma(s):=\gamma_*(s)$ for $s\le s_H$,
$\gamma(s):=\gamma_*(s_H)+(7.5-\gamma_*(s_H))(s-s_H)$ on $[s_H,s_H+1]$, and $\gamma(s):=7.5$ for $s\ge s_H+1$. Define
\[
\log c(z):=\max\big(4\ell^2-\log4,\ \gamma(s)\Lambda\big),\qquad\text{that is,}\qquad
c=\max\big(\tfrac14\qHstd^{1/2},\,X_0^{\gamma(s)}\big).
\]
Then, for every $\Lambda\ge10^8$ and every dyadic height $z\ge X_0$:
\begin{enumerate}[label=\textup{(\alph*)}]
\item (Counting.) The probability, on $\Good(<t_z)$, that some grid window $J$ at height $z$ has
  $\operatorname{rank}S(J)<r_0$ and $|S(J)|\ge r_0+\varphi(r_0)$ is at most $L^{-3}$. For such $J$, let $V$ be a heavy
  subspace of least dimension for $S=S(J)$ (Lemma~\ref{lem:heavy}); then $S':=S(J)\cap V$ satisfies $\qCap(S')\ge c(z)$.
  % Approved deviation from v3 (CORRECTIONS-vs-v3.md, #18): v3 states (a) for any S' of rank < r_0 with Cap(S') >= c.
\item (Absorption compatibility.) If $s\le s_H$, then $y(z)<\Xm=X_0-H(X_0)/2$, so $\Abs$ has no element at height $z$. If
  $s>s_H$, condition \eqref{eq:twelvetwo} of Lemma~\ref{lem:lemmaP} holds at $z$.
\item (Regularity.)
  \begin{itemize}
  \item $c$, $M_1$, $H=\max(\qHstd,M_1^2)$ and $B$ are non-decreasing in $z$.
  \item $d\log H/d\log z\le16\ell/L+10^{-7}\le4\cdot10^{-6}$. Here $16\ell/L$ comes from $\qHstd$, and the raise
    contributes at most $10^{-7}$.
    % Approved deviation from v3 (CORRECTIONS-vs-v3.md, #18): v3 has 4.1e-8 here, from a chain valid only for Lambda = Lambda_s.
    Hence $\rho_H\le1+1.2\cdot10^{-5}$, $H(vw)/H(v)\le w^{4\cdot10^{-6}}$ for
    $v\ge X_0$ and $w\ge1$ (the regularity \eqref{eq:RH}), and $H(v)/v$ is decreasing.
  \item $\log H(z)-\log\qHstd(z)\le10^{-6}\log z$; so the raise adds nothing to the lag.
  \item $c=\tfrac14\qHstd^{1/2}$ for $\ell\ge\ell_{\mathrm{std}}:=\sqrt{(7.5\Lambda+\log4)/4}$, which is $1.37\cdot10^4$ at
    $\Lambda=10^8$.
  \end{itemize}
\end{enumerate}
Parts (b) and (c) are deterministic statements about the parameters. Their proofs (Steps 5 and 6 below) use only the
definitions of $\gamma$, $c$, $M_1$, $L_0$, $y$, $H$ and $B$, and not part (a). In particular (F3$'$) and $\rho_H$, which use
(c), do not depend on (a).
\end{theorem}

\begin{proof}
\emph{Step 0 (reduction).} Let $J$ and $S'=S(J)\cap V$ be as in (a), so $\qCap(S')\ge c$ by Lemma~\ref{lem:heavy}(i). By
Lemma~\ref{lem:reduction}(i) and Theorem~\ref{thm:binomial}, $S'$ yields a configuration with $|R|\ge1+j$ rows outside a minimum cover, where $j\ge(\log c-1)/(2a)-3/2$ and
$a:=\log(r_0(D+1))$. By Lemma~\ref{lem:reduction}(iii) it suffices to show $(1+j)\log(1/w)\ge L+10\ell+10$ for a common bound $w$
on the row weights. We use:
\begin{itemize}
\item $r_0\le L/(2\ell)$ and $L/\Lambda\le D\le2L/\Lambda$;
\item $m:=r_0D\le L^2/(\ell\Lambda)$, so $\log2m\le A$. Also $a\le A-\log2+\log(1+1/D)\le A$, since $D\ge1$, and
  $a\le A-\log2+10^{-9}$ when $D\ge10^9$ (as in Step 3);
\item $A\le2\ell$, and $\log c\ge4\ell^2-\log4\ge3.9\ell^2$ (since $\ell\ge\log\Lambda\ge18.42$).
\end{itemize}

\emph{Step 1 ($G$ is decreasing in $\Lambda$).} At fixed $s$, $\partial A/\partial\Lambda=(3-2/\ell)/\Lambda\in(0,3/\Lambda)$.
So
\[
\partial_\Lambda\log G^2\le\frac3\Lambda\Big(\frac1A+\frac1{A+1}\Big)-\frac{1-3/\Lambda}{\Lambda-A-1}<0,
\]
because $A\ge\log(2\Lambda/\ell)\ge16>6$. Every admissible pair has $\Lambda\ge\Lambda_s$, since $z\ge X_0$ means
$s\ge-\log\Lambda$. Hence $\log c\ge\gamma_*(s)\Lambda\ge1.3\,G(\Lambda,s)\Lambda$ for every $s\le s_H$. Moreover $\gamma_*$
is non-decreasing in $s$, since $s\mapsto e^{s/2}$ increases and $\Lambda\mapsto G$ decreases.

\emph{Step 2 (rows of high degree: $s\le s_H+1$).} Recall $Q=c^{1/2}$.
\begin{itemize}
\item The vectors of degree at most $k$ in $m$ coordinates number at most $\binom{m+k}k\le\exp(k(1+\log2m))$. So the elements
  of degree at most $D^*:=\lfloor\log Q/(1+\log2m)\rfloor$ number at most $Q$.
\item If $D^*\ge D$, then $\qCap\ge c$ is impossible, by Lemma~\ref{lem:heavy}(iii) and (V1). Otherwise Lemma~\ref{lem:reduction}(i)
  completes the cover with rows of degree at least $k:=D^*+1\ge\log c/(2(A+1))$. If no such row exists,
  Lemma~\ref{lem:heavy}(iii) gives a contradiction; this includes the case $D^*=D-1$ with $D$ an integer.
\item Lemma~\ref{lem:rowweights}(b) applies, since $\pi\le m\le\Lambda^3e^{14}\ll X_0$. It gives
  $\log(1/w)\ge k(\Lambda-A-1)-1.1$.
\item Therefore, using $(A+1)/\log c\le(2\ell+1)/(3.9\ell^2)\le0.03$,
  \[
  (1+j)\log(1/w)\ge\frac{(\log c)^2(\Lambda-A-1)}{4A(A+1)}(1-0.031)\ge1.69\cdot0.969\cdot\Lambda^2e^s\ge1.63L\ge L+10\ell+10,
  \]
  as soon as $\log c\ge1.3\,G(\Lambda,s)\Lambda$.
\item This holds for $s\le s_H$ by Step 1. On $[s_H,s_H+1]$ it holds because $\gamma(s)\ge\gamma_*(s)$ there: $\Lambda_s=10^8$ on this range, and $G(10^8,\cdot)$
  grows at logarithmic rate at most $\tfrac12+2/A+10^{-7}\le0.532$ (as $A\ge64.1$ there), so
  $\gamma_*(s_H+t)\le\gamma_*(s_H)e^{0.532t}\le\gamma_*(s_H)+0.24t\le\gamma(s_H+t)$ for $t\in[0,1]$.
\end{itemize}

\emph{Step 3 (Dickman rows: $s\ge s_H+1$).} Here $\log c\ge7.5\Lambda$. We have $D-1\ge\Lambda e^s-1\ge10^4$, and for
$\ell\le0.3\Lambda$ the hypotheses of Lemma~\ref{lem:rowweights}(d) hold, since $\pi X_0^{-1}(eD)^{0.9}\le e^{0.87\Lambda-\Lambda+1}$.
So $\log(1/w)\ge g\ge0.9\Lambda e^s(\log\Lambda+s)(1-10^{-9})$. Also $1+j\ge(\log c-1-a)/(2a)\ge(7.5\Lambda-1-a)/(2a)$, with
$a\le a_+:=3\log\Lambda+2s-2.9\le0.6\Lambda$ (as $\ell\le0.3\Lambda$). Hence, since the bound decreases in $a$,
\[
\frac{(1+j)g}{L}\ge0.45\,(1-10^{-9})\Big(7.5-\frac{1+a_+}{\Lambda}\Big)\frac{\log\Lambda+s}{a_+},\qquad
\frac{\log\Lambda+s}{a_+}=\frac13+\frac{s+2.9}{3a_+}.
\]
If $a_+\le10^{-3}\Lambda$, this is at least $0.45\cdot7.498/3\ge1.12$. Otherwise $s+2.9\ge(a_+-3\log\Lambda)/2$, so
$(\log\Lambda+s)/a_+\ge\tfrac12-10^{-4}$, and the bound is at least $0.45\cdot6.89\cdot0.4999\ge1.5$. In both cases it
exceeds $1+(10\ell+10)/L$.
Where the standard term dominates $\log c$, $j$ is only larger.

\emph{Step 4 (tails).}
\begin{itemize}
\item (T1) $\ell_{\mathrm{std}}\le\ell\le0.3\Lambda$. Here $\log c=4\ell^2-\log4$ and $j\ge\ell-2$. Lemma~\ref{lem:rowweights}(d)
  gives $(1+j)g/L\ge0.9(\ell-1)(\ell-\log\Lambda-0.01)(1-\Lambda/L)/\Lambda\ge0.9\cdot1.87>1.6$.
\item (T2) $\ell\ge0.3\Lambda$. Lemma~\ref{lem:rowweights}(e) with $\eta=\tfrac14$ gives
  $\log(1/w)\ge L/4-5.3X_0^{1/4}/\Lambda-5.6(X_0+L^2)^{1/4}-\log\Lambda-3\ge0.24L$, since $L\ge e^{0.3\Lambda}$. Then $j\ge4$
  suffices.
\end{itemize}
So (a) holds at every scale, including the scales $L\ge X_0^2$ of the zeros.

\emph{Step 5 (proof of (b)).}
\begin{itemize}
\item \emph{$s\le s_H$.} The bounds $\log M_1\le\log c+\log2$, $\log H=\max(8\ell^2,2\log M_1)$ and
  $\log L_0\le\log48-2\log c'+\log M_1+3\ell=7.67+\log M_1+3\ell$ (with $c'=0.15$) give $\log y\le3\log c+9\ell+11.14$. The
  constant is $11.14$ when $\qHstd$ dominates ($8\ell^2\le2\log c+2\log4$), and $9.8$ when $M_1^2$ does. If
  $\log c=\gamma\Lambda\le(\tfrac13-10^{-5})\Lambda$, this is at most $\Lambda-2597<\log\Xm$, since
  $9\ell+11.14\le3\cdot10^{-5}\Lambda-2597$ for $\ell\le2\log\Lambda+6$, and $\log\Xm=\Lambda-O(e^{2715-\Lambda})$. If
  $\log c=4\ell^2-\log4$, it is at most $12\ell^2+9\ell+7\le\Lambda$.
\item \emph{$s\ge s_H$ with $\log c=\gamma\Lambda\le7.5\Lambda$.} Here $\log2y\le3.01\gamma\Lambda$ and
  $\log C_d+4\ell\le1.01\gamma\Lambda$. With $u_0\ge\theta_b\Lambda e^s/(3.01\gamma)$, condition \eqref{eq:twelvetwo} follows from
  $e^s\log(\theta_b\Lambda e^s/(3.01\gamma))\ge107.6\gamma^2$. The left side is increasing in $\Lambda$ and $s$. At
  $\Lambda=10^8$ and $s=s_H+t$ it is at least $393e^t(19.2+t)$, against $107.6(0.334+7.17t)^2\le6053$ on the ramp and
  $=6053$ after it. Also $u_0\ge100$.
\item \emph{Standard regime, $\ell\ge\ell_{\mathrm{std}}$.} Here $\log2y\le12.1\ell^2$, the right side is at most
  $16.3\ell^2$, and $u_0\ge e^\ell/(97\ell^2)$. So \eqref{eq:twelvetwo} holds with an enormous margin.
\end{itemize}

\emph{Step 6 (proof of (c)).}
\begin{itemize}
\item $\gamma$ is non-decreasing, so $c$, $M_1=r_0+\varphi(r_0)+1$, $H$ and $B$ are non-decreasing.
\item $d\log H/dL\le2\Lambda\gamma'(s)/L+16\ell/L$. On the ramp this is at most $2\cdot7.2/(\Lambda e^{s_H})\le4\cdot10^{-10}$.
\item For $s\le s_H$, where $\gamma_*\Lambda\ge4\ell^2-\log4$ (so that $\gamma_*$ dominates the standard term), the raise
  contributes $2\Lambda\gamma_*'(s)/L=2\kappa\gamma_*\Lambda/L$, with $\kappa:=d\log\gamma_*/ds$. Put
  $\ell_s:=2\log\Lambda_s+s$ and $A_s:=A(\Lambda_s,s)$; then $18.42\le\ell_s\le\ell$, since $\Lambda_s\le\Lambda$ and
  $s\ge-\log\Lambda$. Dominance gives $1/\Lambda\le\gamma_*/(4\ell^2-\log4)$, so
  \[
  \frac{\gamma_*\Lambda}L=\frac{\gamma_*}{\Lambda e^s}\le\frac{\gamma_*^2}{(4\ell^2-\log4)e^s}
  =\frac{6.76A_s(A_s+1)}{(4\ell^2-\log4)(\Lambda_s-A_s-1)}\le\frac{1.7A_s(A_s+1)}{\ell_s^2\Lambda_s}.
  \]
  \begin{itemize}
  \item If $s\ge-\log10^8$, then $\Lambda_s=10^8$ and $\ell_s\le42.82$. Here
    $\kappa=\tfrac12+\tfrac12A'(\tfrac1A+\tfrac1{A+1})+\tfrac12A'/(\Lambda_s-A-1)$ with $A'=2-1/\ell_s$, so
    $\kappa\le\tfrac12+2/A_s+10^{-7}\le0.63$, as $A_s\ge16.2$. Also $A_s\le2\ell_s-20.6$, so
    $A_s(A_s+1)/\ell_s^2\le(2-19.6/42.82)^2\le2.38$. The contribution is at most
    $2\cdot0.63\cdot1.7\cdot2.38\cdot10^{-8}\le5.1\cdot10^{-8}$.
  \item If $s<-\log10^8$, then $\Lambda_s=e^{-s}=e^{\ell_s}$ and $A_s=\log2+\ell_s-\log\ell_s$, so $A_s+1\le\ell_s$. The
    variation of $\Lambda_s$ adds at most $\tfrac12(1+10^{-6})$ to the logarithmic rate, while the $A$-terms are now
    non-positive, so $\kappa\le1+10^{-6}$. The contribution is at most $2(1+10^{-6})\cdot1.7\cdot e^{-\ell_s}\le3.5\cdot10^{-8}$.
  \end{itemize}
  So the raise contributes at most $10^{-7}$.
\item $(\log H-\log\qHstd)/L\le\max(0,2\log M_1-8\ell^2)/L\le(2\gamma\Lambda+2)/L\le10^{-6}$ at every scale. \qedhere
\end{itemize}
\end{proof}

\begin{remark}[Numerical cross-check]\label{rem:caplw-numerics}
% v3: §4.6, numerical cross-check after Theorem 4.17
The proof of Theorem~\ref{thm:caplw} is analytic. As an independent check, the precise criteria behind Steps 2--5 were evaluated
numerically with $c$ as above: $D^*$ from $\binom{m+k}k$, the exact $j$, Lemma~\ref{lem:rowweights}(b) and (d), and
\eqref{eq:twelvetwo} with the exact $y$. On the grid $\Lambda\in\{10^8,10^9,10^{10},10^{12},10^{16},10^{20}\}$,
$s\in[-\log\Lambda,30]$ with step $0.01$ (about $35\,000$ points), there are no failures. An earlier cell-by-cell check at
$\Lambda=10^8$ alone, with a knot-based ramp, also passes: $419\,591$ cells, each checked at its worst corners with the degree
cap $D$ evaluated at the top of the cell, and no failures. The maximum of $(\log H-\log\qHstd)/L$ over
$\Lambda\in[10^8,10^{30}]$ is $1.3\cdot10^{-8}$, attained near $s\approx-9.9$ at $\Lambda=10^8$. (Numerical; see
Appendix~\ref{app:cert-capl}.)
\end{remark}

\begin{remark}\label{rem:caplw-sizes}
% v3: §4.6, remarks after the cross-check
At $\Lambda=10^8$, $s_H$ corresponds to $\ell=42.82$ and the ramp ends at $\ell=43.82$. On the raised range,
$c\le X_0^{7.5}$ and the absorption threshold satisfies $y\le X_0^{22.6}L^9$; elsewhere $y=\qHstd L_0L^6=e^{O(\ell^2)}$. In
both cases $y\le n^{o(1)}$, as Lemma~\ref{lem:ZA} requires.
\end{remark}

\subsection{Absorption}\label{sec:queue-abs}
% v3: §4.7 (Lemma P, S1, S2 and the raised threshold)

Recall from Section~\ref{sec:construction} the absorption set
\[
\Abs=\{n\in\NN:\ \log n\ge L_*,\ \text{the latest prime $r$ of $n$ is jittered, and }r<y(n)\},
\]
with $L_*=10^7$ and $y(n):=H(n)L_0(n)(\log n)^6$, the parameters being those of the scale of $n$. Put $\theta_b:=\theta/2$
(so $\theta_b=\tfrac18$ for $\theta=\tfrac14$), $h(u):=u^{\theta_b}$ and $\rho_{\Abs}(u):=\#(\Abs\cap[u,u+h])/h$. The target $T$
gets a unit jump at each $n\in\Abs$, and $T'$ gets $-\rho_{\Abs}$ (see \eqref{eq:target}). Membership is decided when $r$ is
decided, at slot approximately $r<y(n)\ll n$, so $T$ is adapted.

Every jittered prime is at least $\Xm=X_0-H(X_0)/2$ (centred placement), so \emph{$\Abs$ has no element at heights where
$y<\Xm$.} By Theorem~\ref{thm:caplw}(b) these are all heights with $s=\log\log n-2\log\Lambda\le s_H=5.9745$ (that is,
$\log\log n\le42.82$ at $\Lambda=10^8$). So below that height Lemma~\ref{lem:lemmaP} is vacuous, and, since $\Lambda\ge10^8>L_*$, the condition $\log n\ge L_*$
plays no role.

\begin{lemma}[Pointwise sparsity of $\Abs$]\label{lem:lemmaP}
% v3: Lemma 4.18 (Lemma P)
Let $u\ge X_0$ and $L:=\log u$, and put $y:=\sup_{v\in[u,2u]}y(v)$ and $C_d:=M_1(u)L^6/2$. Suppose
\begin{equation}\label{eq:twelvetwo}
0.9u_0\log u_0\ge\log(2y)+\log C_d+4\log L,\qquad u_0:=\frac{\theta_bL}{\log(2y)},\qquad u_0\ge100.
\end{equation}
Then $\#(\Abs\cap[u,u+h])\le h\,L^{-4}$ on $\Good(<t)$ for every $t\ge t_u:=9u/L^{11}+H(u)$. In particular this holds on
$\Good(<t)$ at heights up to $t+1.01H(t)$, since $9(t+1.01H)/L^{11}+H<t$. Condition \eqref{eq:twelvetwo} holds at every $u$ at
which $\Abs$ meets $[u,2u]$, by Theorem~\ref{thm:caplw}(b).
\end{lemma}

\begin{proof}
(i) \emph{Smoothness.} Every prime of $n\in\Abs\cap[u,u+h]$ is decided no later than its latest prime $r<y$, so it has value
less than $r+H(r)<2y$.

(ii) \emph{Divisor trick.} Multiply the primes of $n$ in any order and stop at the first partial product $d>h$. Then
$d\in(h,2yh]$ is a $2y$-smooth element of $\NN$, and $n=d\cdot m$ with $m\in\NN$.

(iii) Hence $\#(\Abs\cap[u,u+h])\le\sum_{d\in\mathcal G_{2y}\cap(h,2yh]}N(I_d)$, where $\mathcal G_{2y}$ is the set of
$2y$-smooth elements of $\NN$ and $I_d:=[u/d,(u+h)/d]$ has length less than $1$.

(iv) \emph{Window count.} $I_d$ lies at height $u/d\le u^{1-\theta_b}$, and it is covered by at most $L^6/4+1$ grid windows
there. On $\Good(<u)\subset\qEcl^{(u)}$ each of them holds at most $M_1(u/d)\le M_1(u)$ elements, so $N(I_d)\le C_d$.
(Alternatively, use Corollary~\ref{cor:localcount} at height $u/d$.)

(v) \emph{Rankin's bound for the realised system.} The primes below $2y$ are the rational primes below $X_0$, together with
jittered primes, which satisfy $\pi_{\NN}(s)\le2s/\log s$ (Lemma~\ref{lem:cm}(i); the local prime number theorem itself is not
needed below $X_{\mathrm{LP}}$, \eqref{eq:XLP}). For $Y\ge h$ put $u_*:=\log Y/\log(2y)$ and $\sigma:=1-\log u_*/\log(2y)$, and suppose $u_*\ge2\Lambda$ and
$\log\log2y\le1.79u_*$. (In use $u_*\ge u_0$, and
$u_0\ge2\Lambda$ wherever $\Abs\ne\emptyset$: by Step 5 of the proof of Theorem~\ref{thm:caplw}, $u_0\ge\theta_b\Lambda e^s/(3.01\gamma)\ge5\Lambda$
on the raised range, and $u_0\ge e^\ell/(97\ell^2)$ beyond it. The same step gives $\log2y\le\max(22.6\Lambda,12.1\ell^2)$,
so $\log\log2y\le1.79u_*$ as well.) Rankin's bound and partial summation as in
Lemma~\ref{lem:rowweights}(c), with $2s/\log s$ in place of $1.26s/\log s$, give
\[
\Psi_{\NN}(Y,2y)/Y\le\exp\big(-u_*\log u_*+2.3u_*/\log u_*+\log\log2y\big)\le\exp(-0.9u_*\log u_*);
\]
here $2.3$ bounds $2(1+2/\log u_*)$ plus the absolute constants, since $\log u_*\ge19$.
For $Y\ge h$ we have $u_*\ge u_0$.

(vi) So $\#(\Abs\cap[u,u+h])\le C_d\Psi_{\NN}(2yh,2y)\le C_d\cdot2yh\cdot\exp(-0.9u_0\log u_0)\le hL^{-4}$, by
\eqref{eq:twelvetwo}.
\end{proof}

No Diophantine input is used: rigid runs are tamed in step (iv), where every count is moved into windows of length less
than $1$.

\begin{lemma}[Consequences for $T'$]\label{lem:rhoA}
% v3: Lemma 4.19
\begin{enumerate}[label=\textup{(\alph*)}]
\item $\sup\rho_{\Abs}\le L^{-4}$ at every height.
% Approved deviation from v3 (CORRECTIONS-vs-v3.md, #16): v3's part (b), an oscillation bound via the local count, is
% deleted (unused, and its proof does not bound the absorbed count). Parts (c), (d) of v3 are (b), (c) here.
\item The source perturbation $\sigma_{\Abs}:=K_{\Abs}\tilde\rho-\rho_{\Abs}$ of Section~\ref{sec:noise} satisfies
  $|\sigma_{\Abs}|\le2L^{-4}$ pointwise. Its compensator form uses $u=L/\log2y\ge10^8$ wherever $\Abs\ne\emptyset$, the raised
  range included.
\item (Current-height form.) On $\Good(<u)$ the $\Abs$-part of $T'$ is $\sigma_{\Abs}$, since $\bar\rho_{\Abs}=\rho_{\Abs}$.
  So the $\Abs$-part of the oscillation term $o_u$ of Section~\ref{sec:noise} satisfies
  \[
  o_{\Abs}(u)\le2\sup|\sigma_{\Abs}|\le4L^{-4}
  \]
  (or $\le2\sup\rho_{\Abs}+2\sup|K_{\Abs}\tilde\rho|\le8L^{-4}$ if the two parts are taken separately). This uses only (a)
  and (b), that is, Lemma~\ref{lem:lemmaP}, which reads $\qEcl$ at heights below $u$. So it involves no queue bound and no event
  $\qLand$ at heights $\ge u$. Then $\sup_v\log v\cdot o_{\Abs}(v)\le8\min L^{-3}\le e^{-120}$ wherever $\Abs\ne\emptyset$,
  since there $L\ge\Lambda^2e^{s_H}$.
\end{enumerate}
Part (c) is the form used where $o_u$ enters R1$'$ at the current height (see \eqref{eq:osc} and Proposition~\ref{prop:pointwise}).
\end{lemma}

\begin{proof}
(a) is Lemma~\ref{lem:lemmaP}. (b) $K_{\Abs}$ is the part of $K_{\mathrm{pt}}$ whose landing is absorbed, that is,
whose latest prime is less than $y(\cdot)$. So its cofactors $w$ are $y$-smooth with $w\ge y/y(y)$. With the \emph{sure} bound
$\tilde\rho\le1.01$ (at most one prime per continuous slot), $K_{\Abs}\tilde\rho\le1.01\sum1/w$ over such $w$. This is at
most $L^{-4}$ by Rankin's bound, with $u=L/\log2y\ge10^8$ and Lemma~\ref{lem:cm}(i) for the realised primes below $y(y)$, which
are decided long before. Together with (a) this gives $|\sigma_{\Abs}|\le2L^{-4}$. No local count at the current height is
used. (c) An oscillation is at most twice the supremum.
\end{proof}

\begin{proposition}[Small latest primes are absorbed]\label{prop:S1S2}
% v3: Proposition 4.20
\begin{enumerate}[label=\textup{(S\arabic*)}]
\item (The small-$q$ truncation in the noise bounds of Lemma~\ref{lem:noise}.) The case of latest primes $q<y$ with jittered
  cofactors needs $\log n<L_*$. Since $L_*=10^7<\Lambda$, it is vacuous. Latest primes
  $q\in[X_0,y)$ are absorbed.
\item (The full blocks of Theorem~\ref{thm:queue} at large heights.) The arrivals that would break \eqref{eq:fullblock} are rigid runs
  $r\cdot r'\cdot d$ driven by small jittered latest primes. These have latest prime less than $y$, so they lie in $\Abs$.
  Non-absorbed arrivals have latest prime $q\ge y$, so Corollary~\ref{cor:increment} applies to them, and \eqref{eq:fullblock} holds as
  shown in Theorem~\ref{thm:queue}.
\end{enumerate}
In particular no small-latest-prime event appears in the budget (Definition~\ref{def:good}).
\end{proposition}

\begin{proof}
This is contained in the statement: (S2) combines the definition of $\Abs$ with Corollary~\ref{cor:increment} and the proof of
Theorem~\ref{thm:queue}(ii); (S1) is the comparison $L_*<\Lambda$.
\end{proof}

\begin{lemma}[Raised threshold]\label{lem:raisedy}
% v3: Lemma 4.21
On the raised range of Theorem~\ref{thm:caplw}, where $\gamma(s)\Lambda>4\ell^2-\log4$, we have
$y(n)\ge y_{\mathrm{std}}(n):=\qHstd(n)L_{0,\mathrm{std}}(n)(\log n)^6$, since $H$, $M_1$ and $L_0$ are all at least their
standard values. So $\Abs_{\mathrm{std}}\subset\Abs$: raising $y$ only \emph{adds} to $\Abs$ the composites whose latest
jittered prime $q$ lies in $[y_{\mathrm{std}}(n),y(n))$. The consumers of $\Abs$ are unaffected:
\begin{itemize}
\item Theorem~\ref{thm:queue} and the noise bounds (Lemma~\ref{lem:noise}) see only latest primes $q\ge y$. Corollary~\ref{cor:increment} needs
  $|I/q|\le H/y=1/(L_0L^6)$, which holds because $y$ uses the same raised $H$ and $L_0$. No consumer assumes $y=y_{\mathrm{std}}$.
  Latest primes in $[y_{\mathrm{std}},y)$ are absorbed, not repaired.
\item Lemma~\ref{lem:S}: its ``$y$'' is the start of the operator, not the absorption threshold. Its bound uses Mertens' bound
  over all of $\NN$, so it is independent of $y$. Its threshold $H(v)\ge(\log v)^{8.5+\delta}$ only gets easier as $H$ grows.
\item Proposition~\ref{prop:pointwise}: its only $\Abs$-input is Lemma~\ref{lem:rhoA}(c), which does not involve $y$ and reads only heights
  below $u$.
\item Lemma~\ref{lem:lemmaP}: \eqref{eq:twelvetwo} holds by Theorem~\ref{thm:caplw}(b). It is read with $y:=\sup_{[u,2u]}y$; along the ramp
  $d\log y/d\log u\le10^{-5}$, so this changes nothing numerically.
\item Lemma~\ref{lem:ZA}: the bound on the absorbed part holds for $y=X_0^{O(1)}\le n^{o(1)}$.
\item Queue size and \eqref{eq:La}: $B\le5M_1L^2/c'\le X_0^{7.6}L^2$ enters additively or as $x^{\eta}$, and is $\ll x^{\theta}$
  at these heights.
\item Causality: $y(n)$ is a deterministic function of the scale of $n$, so membership in $\Abs$ is decided at the decision
  of the latest prime.
\end{itemize}
\end{lemma}

\begin{proof}
Each item is a direct check against the cited statement, as indicated.
\end{proof}

\subsection{Dependencies and non-circularity}\label{sec:queue-deps}
% v3: §4.8

Table~\ref{tab:queuedeps} lists the inputs of the results of this section and of Appendix~\ref{app:rankfull}.

\begin{table}[ht]
\centering\small
\renewcommand{\arraystretch}{1.15}
\begin{tabular}{@{}p{0.3\textwidth}p{0.33\textwidth}p{0.31\textwidth}@{}}
\hline
result & inputs & used in\\
\hline
(F1)--(F2) & Section~\ref{sec:construction} (box jitter, slots) & throughout\\
(F3$'$), (F4) & Lemma~\ref{lem:cm}(iv); Theorem~\ref{thm:caplw}(c); \cite{RosserSchoenfeld1962} & Lemma~\ref{lem:landing}, Lemma~\ref{lem:pil},
  Lemma~\ref{lem:fac}\\
Lemma~\ref{lem:landing} & (F1)--(F2); Theorem~\ref{thm:caplw}(c) (for $\rho_H$); \cite{Freedman1975} & Corollary~\ref{cor:localcount}; event E4\\
Corollary~\ref{cor:localcount} & Lemma~\ref{lem:landing}; Theorem~\ref{thm:queue} at slots $\le\max J+H$; Lemma~\ref{lem:count} &
  Lemma~\ref{lem:threshold-mass}, Lemma~\ref{lem:tau-near}, Proposition~\ref{prop:br}\\
Theorem~\ref{thm:queue} & Corollary~\ref{cor:r1}, Proposition~\ref{prop:pointwise}, Proposition~\ref{prop:ecl}, Corollary~\ref{cor:increment}, (F1)--(F3$'$),
  Lemma~\ref{lem:lindley}, Lemma~\ref{lem:excess}; \cite{Freedman1975} & \eqref{eq:La}, Section~\ref{sec:zeros}, Section~\ref{sec:noise},
  Lemma~\ref{lem:budget}\\
Definition~\ref{def:good}, Lemma~\ref{lem:budget} & Theorem~\ref{thm:queue}, Proposition~\ref{prop:ecl}, Lemma~\ref{lem:landing}, Lemma~\ref{lem:noise} & Theorem~\ref{thm:A}\\
Proposition~\ref{prop:ecl} & Theorem~\ref{thm:R}, Lemma~\ref{lem:heavy}, Theorem~\ref{thm:caplw} & Corollary~\ref{cor:increment}, Theorem~\ref{thm:queue},
  Lemma~\ref{lem:lemmaP}\\
Corollary~\ref{cor:increment} & Proposition~\ref{prop:ecl} & Theorem~\ref{thm:queue}, Lemma~\ref{lem:noise}, Proposition~\ref{prop:pointwise}\\
Lemma~\ref{lem:heavy} & none & Proposition~\ref{prop:ecl}, Lemma~\ref{lem:reduction}, Theorem~\ref{thm:caplw}\\
Lemma~\ref{lem:covering}, Lemma~\ref{lem:coefsupport}, Theorem~\ref{thm:binomial}, Theorem~\ref{thm:tcl} & none (deterministic) &
  Lemma~\ref{lem:reduction}, Theorem~\ref{thm:caplw}\\
Proposition~\ref{prop:cover} & Lemma~\ref{lem:volume}, Lemma~\ref{lem:fac}, Theorem~\ref{thm:svl}, Lemma~\ref{lem:SC}, (V1) & Lemma~\ref{lem:reduction}\\
Lemma~\ref{lem:reduction} & Lemma~\ref{lem:heavy}, Lemma~\ref{lem:covering}, Theorem~\ref{thm:binomial}, Proposition~\ref{prop:cover}, (F4) &
  Theorem~\ref{thm:caplw}\\
Lemma~\ref{lem:rowweights} & Rankin's bound; \cite{RosserSchoenfeld1962} & Theorem~\ref{thm:caplw}\\
Theorem~\ref{thm:caplw}(a) & Lemma~\ref{lem:heavy}, Lemma~\ref{lem:reduction}, Lemma~\ref{lem:rowweights}, Theorem~\ref{thm:binomial} & Proposition~\ref{prop:ecl}\\
Theorem~\ref{thm:caplw}(b), (c) & none (definitions of the parameters) & the parameters of Section~\ref{sec:queue-conv}, (F3$'$),
  Lemma~\ref{lem:landing}, Theorem~\ref{thm:queue}, Lemma~\ref{lem:lemmaP}\\
Lemma~\ref{lem:lemmaP} & Proposition~\ref{prop:ecl} (heights $<u$), Lemma~\ref{lem:cm}, Theorem~\ref{thm:caplw}(b) & Lemma~\ref{lem:rhoA}, \eqref{eq:target},
  Lemma~\ref{lem:ZA}, Section~\ref{sec:noise} ($\sigma_{\Abs}$)\\
Lemma~\ref{lem:rhoA} & Lemma~\ref{lem:lemmaP} only (heights $<u$) & (c): \eqref{eq:osc}, Proposition~\ref{prop:pointwise} (current height)\\
Proposition~\ref{prop:S1S2} & Lemma~\ref{lem:lemmaP}, Corollary~\ref{cor:increment} & Theorem~\ref{thm:queue}, Lemma~\ref{lem:noise}\\
Lemma~\ref{lem:raisedy} & Theorem~\ref{thm:caplw}, Lemma~\ref{lem:lemmaP}, Lemma~\ref{lem:rhoA} & Lemma~\ref{lem:S}, Proposition~\ref{prop:pointwise},
  Lemma~\ref{lem:ZA}\\
Lemmas~\ref{lem:volume}--\ref{lem:fac}, Lemma~\ref{lem:SC}, Theorem~\ref{thm:svl} & (F1)--(F4) & Theorem~\ref{thm:R}, Proposition~\ref{prop:cover}\\
Theorem~\ref{thm:R} & Lemmas~\ref{lem:volume}--\ref{lem:fac}, Lemma~\ref{lem:SC}, Theorem~\ref{thm:svl} & Proposition~\ref{prop:ecl}\\
\hline
\end{tabular}
\caption{Dependencies of the results of Section~\ref{sec:queue} and Appendix~\ref{app:rankfull}.}\label{tab:queuedeps}
\end{table}

\paragraph{Non-circularity.}
\begin{itemize}
\item Theorem~\ref{thm:queue} at slot $s_j$ reads only events with times at most $s_j$: $\qEcl$ at cofactor heights at most
  $4u/L^{11}$, and Corollary~\ref{cor:r1} and Proposition~\ref{prop:pointwise}.
\item $\qEcl$ reads no queue bound at its own heights. Its inputs are (F1)--(F4) at heights at most $s-H$, carried by
  progressive events.
\item Corollary~\ref{cor:localcount} is used only at strictly smaller heights.
\item Lemma~\ref{lem:lemmaP} uses $\qEcl$ only at heights $u/d<u$.
\item The first-failure argument (Lemma~\ref{lem:budget}) orders all of these by event time.
\end{itemize}
